
\documentclass[11pt,a4paper]{article}
% Shared preamble for Wolf lecture note transcriptions
% Usage: % Shared preamble for Wolf lecture note transcriptions
% Usage: % Shared preamble for Wolf lecture note transcriptions
% Usage: \input{preamble} in each chapter file
% Include guard: safe to \input multiple times (e.g. from main_wolf.tex)
\ifdefined\wolfpreambleloaded\endinput\fi
\def\wolfpreambleloaded{}

\usepackage{amsmath,amssymb,amsthm}
% Match Wolf book-class layout: a4paper, ~345pt text width
\usepackage[a4paper, textwidth=345pt, textheight=550pt]{geometry}
\usepackage{hyperref}
\usepackage{mathtools}

% ── Notation (consistent across all chapters) ──
\newcommand{\mcM}{\mathcal{M}}       % matrix algebra M_d(C)
\newcommand{\mcB}{\mathcal{B}}       % bounded operators B(H)
\newcommand{\mcA}{\mathcal{A}}       % multiplicative domain
\newcommand{\mcH}{\mathcal{H}}       % Hilbert space H
\newcommand{\diag}{\operatorname{diag}}
\newcommand{\one}{\mathbb{1}}        % identity operator
\newcommand{\CC}{\mathbb{C}}
\newcommand{\RR}{\mathbb{R}}
\newcommand{\NN}{\mathbb{N}}
\newcommand{\ZZ}{\mathbb{Z}}
\newcommand{\tr}{\operatorname{tr}}
\newcommand{\spec}{\operatorname{spec}}
\newcommand{\rank}{\operatorname{rank}}
\newcommand{\supp}{\operatorname{supp}}
\newcommand{\range}{\operatorname{range}}
\newcommand{\spn}{\operatorname{span}}
\newcommand{\FT}{\mathcal{F}_T}      % fixed point set
\newcommand{\XT}{\mathcal{X}_T}      % peripheral eigenspace
\newcommand{\mcX}{\mathcal{X}}       % peripheral eigenspace (bare)
\newcommand{\id}{\mathrm{id}}        % identity map
\newcommand{\ket}[1]{|#1\rangle}
\newcommand{\bra}[1]{\langle#1|}
\newcommand{\braket}[2]{\langle#1|#2\rangle}
\newcommand{\ketbra}[2]{|#1\rangle\!\langle#2|}

% ── Helper: properly undefine a theorem environment for redefinition ──
% Usage: \UndefEnv{theorem} before \newtheorem{theorem}{...}
\makeatletter
\newcommand{\UndefEnv}[1]{%
  % Remove from section reset list (if registered there)
  \@removefromreset{#1}{section}%
  \expandafter\let\csname #1\endcsname\relax
  \expandafter\let\csname end#1\endcsname\relax
  \expandafter\let\csname c@#1\endcsname\relax
}
\makeatother

% ── Theorem environments ──
\theoremstyle{plain}
\newtheorem{theorem}{Theorem}[section]
\newtheorem{proposition}[theorem]{Proposition}
\newtheorem{corollary}[theorem]{Corollary}
\newtheorem{lemma}[theorem]{Lemma}
\theoremstyle{definition}
\newtheorem{definition}[theorem]{Definition}
\newtheorem{example}[theorem]{Example}
\newtheorem{problem}[theorem]{Problem}
\theoremstyle{remark}
\newtheorem{remark}[theorem]{Remark} in each chapter file
% Include guard: safe to \input multiple times (e.g. from main_wolf.tex)
\ifdefined\wolfpreambleloaded\endinput\fi
\def\wolfpreambleloaded{}

\usepackage{amsmath,amssymb,amsthm}
% Match Wolf book-class layout: a4paper, ~345pt text width
\usepackage[a4paper, textwidth=345pt, textheight=550pt]{geometry}
\usepackage{hyperref}
\usepackage{mathtools}

% ── Notation (consistent across all chapters) ──
\newcommand{\mcM}{\mathcal{M}}       % matrix algebra M_d(C)
\newcommand{\mcB}{\mathcal{B}}       % bounded operators B(H)
\newcommand{\mcA}{\mathcal{A}}       % multiplicative domain
\newcommand{\mcH}{\mathcal{H}}       % Hilbert space H
\newcommand{\diag}{\operatorname{diag}}
\newcommand{\one}{\mathbb{1}}        % identity operator
\newcommand{\CC}{\mathbb{C}}
\newcommand{\RR}{\mathbb{R}}
\newcommand{\NN}{\mathbb{N}}
\newcommand{\ZZ}{\mathbb{Z}}
\newcommand{\tr}{\operatorname{tr}}
\newcommand{\spec}{\operatorname{spec}}
\newcommand{\rank}{\operatorname{rank}}
\newcommand{\supp}{\operatorname{supp}}
\newcommand{\range}{\operatorname{range}}
\newcommand{\spn}{\operatorname{span}}
\newcommand{\FT}{\mathcal{F}_T}      % fixed point set
\newcommand{\XT}{\mathcal{X}_T}      % peripheral eigenspace
\newcommand{\mcX}{\mathcal{X}}       % peripheral eigenspace (bare)
\newcommand{\id}{\mathrm{id}}        % identity map
\newcommand{\ket}[1]{|#1\rangle}
\newcommand{\bra}[1]{\langle#1|}
\newcommand{\braket}[2]{\langle#1|#2\rangle}
\newcommand{\ketbra}[2]{|#1\rangle\!\langle#2|}

% ── Helper: properly undefine a theorem environment for redefinition ──
% Usage: \UndefEnv{theorem} before \newtheorem{theorem}{...}
\makeatletter
\newcommand{\UndefEnv}[1]{%
  % Remove from section reset list (if registered there)
  \@removefromreset{#1}{section}%
  \expandafter\let\csname #1\endcsname\relax
  \expandafter\let\csname end#1\endcsname\relax
  \expandafter\let\csname c@#1\endcsname\relax
}
\makeatother

% ── Theorem environments ──
\theoremstyle{plain}
\newtheorem{theorem}{Theorem}[section]
\newtheorem{proposition}[theorem]{Proposition}
\newtheorem{corollary}[theorem]{Corollary}
\newtheorem{lemma}[theorem]{Lemma}
\theoremstyle{definition}
\newtheorem{definition}[theorem]{Definition}
\newtheorem{example}[theorem]{Example}
\newtheorem{problem}[theorem]{Problem}
\theoremstyle{remark}
\newtheorem{remark}[theorem]{Remark} in each chapter file
% Include guard: safe to \input multiple times (e.g. from main_wolf.tex)
\ifdefined\wolfpreambleloaded\endinput\fi
\def\wolfpreambleloaded{}

\usepackage{amsmath,amssymb,amsthm}
% Match Wolf book-class layout: a4paper, ~345pt text width
\usepackage[a4paper, textwidth=345pt, textheight=550pt]{geometry}
\usepackage{hyperref}
\usepackage{mathtools}

% ── Notation (consistent across all chapters) ──
\newcommand{\mcM}{\mathcal{M}}       % matrix algebra M_d(C)
\newcommand{\mcB}{\mathcal{B}}       % bounded operators B(H)
\newcommand{\mcA}{\mathcal{A}}       % multiplicative domain
\newcommand{\mcH}{\mathcal{H}}       % Hilbert space H
\newcommand{\diag}{\operatorname{diag}}
\newcommand{\one}{\mathbb{1}}        % identity operator
\newcommand{\CC}{\mathbb{C}}
\newcommand{\RR}{\mathbb{R}}
\newcommand{\NN}{\mathbb{N}}
\newcommand{\ZZ}{\mathbb{Z}}
\newcommand{\tr}{\operatorname{tr}}
\newcommand{\spec}{\operatorname{spec}}
\newcommand{\rank}{\operatorname{rank}}
\newcommand{\supp}{\operatorname{supp}}
\newcommand{\range}{\operatorname{range}}
\newcommand{\spn}{\operatorname{span}}
\newcommand{\FT}{\mathcal{F}_T}      % fixed point set
\newcommand{\XT}{\mathcal{X}_T}      % peripheral eigenspace
\newcommand{\mcX}{\mathcal{X}}       % peripheral eigenspace (bare)
\newcommand{\id}{\mathrm{id}}        % identity map
\newcommand{\ket}[1]{|#1\rangle}
\newcommand{\bra}[1]{\langle#1|}
\newcommand{\braket}[2]{\langle#1|#2\rangle}
\newcommand{\ketbra}[2]{|#1\rangle\!\langle#2|}

% ── Helper: properly undefine a theorem environment for redefinition ──
% Usage: \UndefEnv{theorem} before \newtheorem{theorem}{...}
\makeatletter
\newcommand{\UndefEnv}[1]{%
  % Remove from section reset list (if registered there)
  \@removefromreset{#1}{section}%
  \expandafter\let\csname #1\endcsname\relax
  \expandafter\let\csname end#1\endcsname\relax
  \expandafter\let\csname c@#1\endcsname\relax
}
\makeatother

% ── Theorem environments ──
\theoremstyle{plain}
\newtheorem{theorem}{Theorem}[section]
\newtheorem{proposition}[theorem]{Proposition}
\newtheorem{corollary}[theorem]{Corollary}
\newtheorem{lemma}[theorem]{Lemma}
\theoremstyle{definition}
\newtheorem{definition}[theorem]{Definition}
\newtheorem{example}[theorem]{Example}
\newtheorem{problem}[theorem]{Problem}
\theoremstyle{remark}
\newtheorem{remark}[theorem]{Remark}

% --- Wolf-style numbering: each environment has its own counter (within Chapter/Section 6) ---
% The shared preamble defines proposition/lemma/etc. sharing the theorem counter;
% Wolf numbers them independently (Prop 6.1 and Thm 6.1 both occur), so we override locally.
\UndefEnv{proposition}
\UndefEnv{corollary}
\UndefEnv{lemma}
\UndefEnv{definition}
\UndefEnv{example}
\UndefEnv{problem}
\UndefEnv{remark}

\theoremstyle{plain}
\newtheorem{proposition}{Proposition}[section]
\newtheorem{corollary}{Corollary}[section]
\newtheorem{lemma}{Lemma}[section]
\theoremstyle{definition}
\newtheorem{definition}{Definition}[section]
\newtheorem{example}{Example}[section]
\newtheorem{problem}{Problem} % global numbering (Problem 12, 13, ...)
\theoremstyle{remark}
\newtheorem{remark}{Remark}[section]

% \mcH, \diag now in preamble

\begin{document}

% Standalone compilation note:
% We set the section counter so that the following \section is numbered "6".
% Remove/adjust this if you integrate into a master document.
\setcounter{section}{5}

\section{Spectral properties}

If a linear map has equal input and output space, like
\[
    T : \mcM_d(\CC)\to \mcM_d(\CC),
\]
we can assign a spectrum to it. As outlined in Sec.~1.6 the spectrum can be
defined in purely algebraic terms as the set of complex numbers $\lambda$ for
which $(\lambda \id-T)$ is not invertible. Since we restrict ourselves to finite
dimensions this is equivalent to the set of $\lambda$'s for which there is an
$X\in \mcM_d(\CC)$ such that
\begin{equation}\tag{6.1}\label{eq:6.1}
    T(X)=\lambda X.
\end{equation}
Although $X$ is an operator, we will refer to it as an eigenvector of $T$ as
$\mcM_d(\CC)$ is considered as a vector space in this context. Hence, the
spectrum of $T$ is the collection of eigenvalues of the $d^2\times d^2$ matrix
representation $\widehat{T}$ introduced in Sec.~2.3. Using this representation
we see that the eigenvalues of $T$ and $T^*$ coincide due to
$\widehat{T^*}=\widehat{T}^{\mathsf T}$ when using a Hermitian basis of operators
in Eq.~(2.20).

\paragraph{Location of eigenvalues.}
For Hermiticity preserving maps (i.e., in particular for positive maps) the
adjoint of Eq.~\eqref{eq:6.1} implies that eigenvalues are either real or they
come in complex conjugate pairs.\footnote{More precisely, if we decompose $X$
    into Hermitian and anti-Hermitian parts and use the Hermiticity preservation
    together with linearity we obtain $T(X^\dagger)=\overline{\lambda}\,X^\dagger$.}
Moreover, for every real $\lambda$ there is at least one corresponding Hermitian
eigenvector $X$.

Recall that the spectral radius is defined as
\[
    \varrho(T):=\sup\{|\lambda| \mid \lambda\in\spec(T)\},
\]
i.e., as the eigenvalue which is largest in magnitude.

\begin{proposition}[Spectral radius of positive maps]
    If $T$ is a positive map on $\mcM_d(\CC)$, then its spectral radius satisfies
    \begin{equation}\tag{6.2}\label{eq:6.2}
        \varrho(T)\le \|T(\one)\|_\infty.
    \end{equation}
    If in addition $T$ is unital or trace-preserving, then there is an eigenvalue
    $\lambda=1$. That is, in this case $\varrho(T)=1$ so that all eigenvalues of $T$
    lie in the unit disc of the complex plane.
\end{proposition}

\begin{proof}
    By the theorem of Russo and Dye we have $\|T(X)\|_\infty \le \|T(\one)\|_\infty\|X\|_\infty$.
    Therefore if $T(X)=\lambda X$, then
    \begin{equation}\tag{6.3}\label{eq:6.3}
        |\lambda|\|X\|_\infty = \|T(X)\|_\infty \le \|T(\one)\|_\infty\|X\|_\infty,
    \end{equation}
    which implies Eq.~\eqref{eq:6.2}. If $T$ is unital then $T(\one)=\one$ provides
    an eigenvalue $\lambda=1$. Similar holds if $T$ is trace-preserving, since then
    $T^*$ (which has the same spectrum) is unital.
\end{proof}

Another way to see that the spectral radius of every quantum channel (or
trace-preserving positive map) is one would be to use that
$\varrho(\widehat{T})=\lim_{n\to\infty}\|\widehat{T}^n\|^{1/n}$ (see Eq.~(1.35))
together with the fact that $T^n$ remains a quantum channel for every $n$ and
that quantum channels in finite dimensions have bounded norm (see \dots).

To summarize, the eigenvalues of every quantum channel are restricted to lie on
the unit disc, there is one eigenvalue one and the others are real or come in
complex conjugate pairs. In fact, there is no further restriction (as in the
case of `classical' stochastic matrices): for every $d>1$, every $\lambda$ with
$|\lambda|\le 1$ can appear as an eigenvalue of a quantum channel on $\mcM_d$.
Only the set of eigenvalues as a whole underlies further non-trivial constraints
(see Sec.~6.1).

\paragraph{Spectral decomposition.}
Recall some basic linear algebra. Like every other square matrix,
$\widehat{T}\in \mcM_{d^2}(\CC)$ admits a Jordan decomposition of the form
\begin{equation}\tag{6.4}\label{eq:6.4}
    \widehat{T}
    = X\left(\bigoplus_{k=1}^K J_k(\lambda_k)\right)X^{-1},
    \qquad
    J_k(\lambda):=
    \begin{pmatrix}
        \lambda &1\\
                &\ddots &1\\
                &       &\lambda
    \end{pmatrix}
    \in \mcM_{d_k}(\CC),
\end{equation}
where the $J_k$'s are Jordan blocks of size $d_k$ (with $\sum_k d_k=d^2$) and the
number $K$ of Jordan blocks equals the number of different non-zero eigenvectors.
It is sometimes useful to subdivide each Jordan block into a projection%
\footnote{This implies that in Eq.~\eqref{eq:6.5} $P_k^2=P_k$ but not necessarily
    $P_k=P_k^\dagger$, i.e., $P_k$ need not be an orthogonal projection.}
and a nilpotent part so that we get
\begin{equation}\tag{6.5}\label{eq:6.5}
    \begin{aligned}
        \widehat{T}
               & = \sum_{k=1}^K (\lambda_k P_k + N_k),
               & N_k^{d_k}                             & =0,
               & N_kP_k                                & = P_kN_k = N_k, \\
        P_kP_l & = \delta_{kl}P_k,
               & \tr[P_k]                              & = d_k,
               & \sum_k P_k                            & = \one.
    \end{aligned}
\end{equation}
The number of Jordan blocks with eigenvalue $\lambda$ is the geometric
multiplicity of $\lambda$, while their joint dimension $\sum_{k:\lambda_k=\lambda} d_k$
is its algebraic multiplicity. If these two multiplicities are equal for every
eigenvalue, then $\widehat{T}$ is called non-defective. Non-defective matrices
are dense (which is easily seen by perturbing the eigenvalues) and exactly those
for which complete bases of eigenvectors exist. We can then write
\begin{equation}\tag{6.6}\label{eq:6.6}
    \widehat{T}=\sum_k \lambda_k \ket{R_k}\!\bra{L_k},
    \qquad
    \langle L_k, R_l\rangle=\delta_{kl},
\end{equation}
so that the right eigenvectors $\{R_k\}$ and the left eigenvectors $\{L_k\}$ are
biorthogonal.

Returning to the fact that $T$ is a linear map on $\mcM_d(\CC)$ this means that
in the non-defective case there are biorthogonal operator bases
$\{R_k\in\mcM_d(\CC)\}$ and $\{L_k\in\mcM_d(\CC)\}$ such that
\begin{equation}\tag{6.7}\label{eq:6.7}
    T(A)=\sum_k \lambda_k\,\tr[L_k^\dagger A]\,R_k,
    \qquad
    \tr[L_k^\dagger R_l]=\delta_{kl},
\end{equation}
so that
\begin{equation}\tag{6.8}\label{eq:6.8}
    \forall k:\quad T(R_k)=\lambda_k R_k.
\end{equation}
The Choi--Jamiolkowski operator $\tau=(T\otimes \id)(\ket{\Omega}\!\bra{\Omega})$
then takes the form
\begin{equation}\tag{6.9}\label{eq:6.9}
    \tau=\frac{1}{d}\sum_k \lambda_k\,\overline{L_k}\otimes R_k.
\end{equation}
If $T$ is a trace-preserving map, then all right eigenvectors $R_k$ corresponding
to eigenvalues $\lambda_k\neq 1$ have to be traceless, $\tr[R_k]=0$ since $\one$
is a left eigenvector.

Note that if $T$ is a quantum channel, then $\widehat{T}$ is generally neither
Hermitian (unless $T=T^*$, see Prop.~2.6) nor is it necessarily non-defective.
However, the uniformly bounded norm of trace-preserving positive maps prevents
eigenvalues of magnitude one (the so-called peripheral spectrum) from having
non-trivial Jordan blocks (meaning $d_k>1$):

\begin{proposition}[Trivial Jordan blocks for peripheral spectrum]
    Let $T:\mcM_d(\CC)\to\mcM_d(\CC)$ be a trace-preserving (or unital) positive linear
    map. If $\lambda$ is an eigenvalue of $T$ with $|\lambda|=1$, then its geometric
    multiplicity equals its algebraic multiplicity, i.e., all Jordan blocks for
    $\lambda$ are one-dimensional.
\end{proposition}

\begin{proof}
    First note that for every trace-preserving (or unital) positive map $\tr[A\,T(B)]$
    can be upper bounded in terms of the dimension $d$ and the norms of $A$ and $B$
    only. While we will discuss such bounds in greater detail later it suffices for
    now to note that we may write $A$ and $B$ as a linear combination of positive
    matrices and that for positive $A,B$ we have
    \[
        \tr[A\,T(B)] \le \|A\|_\infty\|B\|_\infty\tr[\one\,T(\one)]
        = d\|A\|_\infty\|B\|_\infty.
    \]
    The same bound has to hold for every power $T^n$, $n\in\NN$. However, if we take
    the $n$th power of a Jordan block $J(\lambda)$, it is a Toeplitz matrix whose
    first row reads
    \begin{equation}\tag{6.10}\label{eq:6.10}
        [J(\lambda)^n]_{1j}
        = \lambda^{n-j+1}\binom{n}{j-1},
        \qquad
        j=1,\dots,\dim(J(\lambda)).
    \end{equation}
    Hence, if $\dim(J(\lambda))>1$ and $|\lambda|=1$ this would contain entries which
    grow unboundedly with $n$ contradicting the uniform boundedness of $T^n$.
\end{proof}

We can now use the obtained results in order to show that some of the spectral
projections appearing in Eq.~\eqref{eq:6.5} are (completely) positive in their
own right. Define
\begin{equation}\tag{6.11}\label{eq:6.11}
    \widehat{T}_\infty := \sum_{k:\lambda_k=1} P_k,
\end{equation}
\begin{equation}\tag{6.12}\label{eq:6.12}
    \widehat{T}_\phi := \sum_{k:|\lambda_k|=1} P_k,
\end{equation}
\begin{equation}\tag{6.13}\label{eq:6.13}
    \widehat{T}_\varphi := \sum_{k:|\lambda_k|=1} \lambda_k P_k,
\end{equation}
from the spectral decomposition in Eq.~\eqref{eq:6.5} and the corresponding maps
$T_\infty,T_\phi,T_\varphi$ via Eq.~(2.20).

\begin{proposition}[Cesaro means]
    Let $T:\mcM_d(\CC)\to\mcM_d(\CC)$ be a linear map which is trace-preserving and
    (completely) positive. Then $T_\infty,T_\phi,T_\varphi$ are trace-preserving and
    (completely) positive as well. More precisely, we have that
    \begin{enumerate}
        \item[(i)] there exists an increasing sequence $n_i\in\NN$ such that
              $\lim_{i\to\infty} T^{n_i}=T_\phi$,
        \item[(ii)] $T_\varphi = T\,T_\phi$ and
        \item[(iii)]
              \begin{equation}\tag{6.14}\label{eq:6.14}
                  T_\infty=\lim_{N\to\infty}\frac{1}{N}\sum_{n=1}^N T^n.
              \end{equation}
    \end{enumerate}
\end{proposition}

\begin{proof}
    Showing (i)--(iii) essentially proves the claim that the three maps are
    trace-preserving and (completely) positive. Throughout we will use that the
    spectral radius of $T$ is one, that there is an eigenvalue one and that the
    peripheral spectrum of $T$ only has one-dimensional Jordan blocks. For (i) we
    exploit the subsequent Lemma on simultaneous Diophantine approximations. For
    our purpose Lemma~6.1 implies that for every $\varepsilon>0$ there is an $n\in\NN$
    such that $|\lambda_k^n-1|\le \varepsilon$ for all eigenvalues $\lambda_k$ of
    magnitude one. This means that there is an increasing sequence $n_i\in\NN$ such
    that $T^{n_i}$ is a better and better approximation to $T_\phi$. Since the set
    of (completely) positive, trace-preserving maps on $\mcM_d$ is compact, the
    limit point $T_\phi$ belongs to this set. (ii) is obvious from the spectral
    decomposition in Eq.~\eqref{eq:6.5} and (iii) is an immediate consequence of the
    geometric series $\sum_{n=1}^N \lambda^n = (\lambda-\lambda^{N+1})/(1-\lambda)$ for
    $\lambda\neq 1$.
\end{proof}

In a similar vein as $T_\infty$ is expressed as a Cesaro mean in Eq.~\eqref{eq:6.14}
we can write
\begin{equation}\tag{6.15}\label{eq:6.15}
    T_\phi
    = \lim_{N\to\infty}\frac{1}{N}\sum_{n=1}^N
    \sum_{k:|\lambda_k|=1} (\overline{\lambda_k}T)^n.
\end{equation}

The maps $T_\infty$ and $T_\phi$ are projections onto the set of fixed points
(see Chap.~6.4) and the space spanned by the eigenvectors of the peripheral
spectrum respectively. That is, $X\in\mcM_d$ is in the range of $T_\infty$ iff
$T(X)=X$ and it is in the range of $T_\phi$ iff for every $\varepsilon>0$ there
is an $n\in\NN$ such that $\|T^n(X)-X\|\le \varepsilon$. In other words the image
of $T_\phi$ is the set of operators for which recurrences arise.

For completeness (and later use) Dirichlet's theorem on simultaneous Diophantine
approximations which was used in the above proof:

\begin{lemma}[Dirichlet's theorem]
    Let $x_1,\dots,x_m$ be $m$ real numbers and $q>1$ any integer. Then there exist
    integers $n,p_1,\dots,p_m$ such that
    \begin{equation}\tag{6.16}\label{eq:6.16}
        1\le n\le q^m
        \qquad\text{and}\qquad
        |x_k n - p_k|\le \frac{1}{q},
        \ \forall k.
    \end{equation}
\end{lemma}

\paragraph{Resolvents.}
The resolvent $R$ of a linear operator $T$ is an operator-valued function on
$\CC$ defined as
\begin{equation}\tag{6.17}\label{eq:6.17}
    R(z):=(z\id - T)^{-1}.
\end{equation}
The singularities of $R$ are exactly the eigenvalues of $T$. Hence, $R$ is
meromorphic and defined on the set which is the complement of the spectrum of
$T$, called resolvent set. For $|z|>\varrho(T)$, i.e., outside the spectral radius
of $T$, we can express the resolvent as
\begin{equation}\tag{6.18}\label{eq:6.18}
    R(z)=\frac{1}{z}\sum_{t=0}^\infty \frac{T^t}{z^t}.
\end{equation}

If $T^t$ describes the discrete time evolution of a quantum (or classical) system
this series implies that the spectrum of $T$ is reflected in a peculiar way in
the time evolution of expectation values: consider any $\rho,A\in\mcM_d(\CC)$ and
the sequence $\langle A(t)\rangle := \tr[\rho\,T^t(A)]$ for $t=0,1,2,\dots$. If $A$
represents an observable and $\rho$ is a density matrix evolving according to $T$,
then $\langle A(t)\rangle$ is the expectation value at time $t$. The $z$-transform
(or discrete Laplace transform) of the `signal' $\langle A(t)\rangle$ is given by
\begin{equation}\tag{6.19}\label{eq:6.19}
    \mathcal{L}(z)
    := \frac{1}{z}\sum_{t\in\NN_0}\frac{\langle A(t)\rangle}{z^t},
\end{equation}
which converges for $|z|>\varrho(T)$ and can be defined in the interior of that
disc by analytic continuation. Obviously, the poles of $\mathcal{L}$ are located
at eigenvalues of $T$. Note, however, that depending on $A$ and $\rho$ not every
eigenvalue of $T$ will give rise to a pole of $\mathcal{L}$.

Using the resolvent, many of the results of analytic function theory carry over
to the world of operators. For instance if $\Delta\subseteq\CC$ is simply
connected and enclosed by a smooth curve $\partial\Delta$, then
\begin{equation}\tag{6.20}\label{eq:6.20}
    \sum_{k:\lambda_k\in\Delta} P_k
    = \frac{1}{2\pi i}\int_{\partial\Delta} R(z)\,dz,
\end{equation}
and
\begin{equation}\tag{6.21}\label{eq:6.21}
    f(T)
    = \frac{1}{2\pi i}\int_{\partial\Delta} f(z)R(z)\,dz,
    \qquad
    \text{if } \lambda_k\in\Delta\ \forall k,
\end{equation}
where $f:\CC\to\CC$ is any function holomorphic in $\Delta$ and the $P_k$'s are the
projections appearing in the spectral decomposition in Eq.~\eqref{eq:6.5}.

% ------------------------------------------------------------
\subsection{Determinants}
% ------------------------------------------------------------

The determinant (=product of all eigenvalues) is a useful tool in particular when
it comes to the composition of maps since
\begin{equation}\tag{6.22}\label{eq:6.22}
    \det(T_1T_2)=\det(T_1)\det(T_2).
\end{equation}
In fact, the property of being multiplicative essentially characterizes the
determinant: if a functional $F:\mcM_D(\CC)\to\CC$ satisfies
$F(T_1T_2)=F(T_1)F(T_2)$ for all $T_i\in\mcM_D(\CC)$, then
$F(T)=f(\det(T))$ for some multiplicative function $f:\CC\to\CC$~[?].

Also note that Eq.~\eqref{eq:6.22} implies the simple relation $\prod_i s_i=|\det(T)|$
between the determinant and the singular values $\{s_i\}$ of a map.

\begin{theorem}[Determinants]
    Let $T:\mcM_d\to\mcM_d$ be a positive and trace preserving linear map.
    \begin{enumerate}
        \item $\det T$ is real and contained in the interval $[-1,1]$,
        \item $|\det T|=1$ iff $T$ is either a unitary conjugation (i.e., $T(A)=UAU^\dagger$
              for some unitary $U\in\mcM_d$) or unitarily equivalent to a matrix transposition
              (i.e., $T(A)=UA^{\mathsf T}U^\dagger$ for some unitary $U\in\mcM_d$),
        \item if $T$ is a unitary conjugation then $\det T=1$ and if $\det T=-1$ then $T$
              is a matrix transposition up to unitary equivalence. In both cases the converse
              holds iff $\lfloor d/2\rfloor$ is odd.
    \end{enumerate}
\end{theorem}

\begin{proof}
    Since eigenvalues of positive maps either come in complex conjugate pairs or are
    real, their product $\det T$ is real, too. Due to the additional trace-preserving
    property the spectral radius is one (Prop.~6.1) so that $\det T\in[-1,1]$.

    Now consider the case $\det T=\pm 1$ where all eigenvalues are phases. Following
    Dirichlet's Theorem (Lem.~6.1) as in the first part of Prop.~6.3, there is always
    a subsequence $n_i$ such that the limit of powers $\lim_{i\to\infty} T^{n_i}=:T_\phi$
    has eigenvalues which all converge to one which implies that $T_\phi=\id$. Hence,
    the inverse $T^{-1} = T_\phi T^{-1} = \lim_{i\to\infty} T^{n_i-1}$ is a trace-preserving
    positive map as well.

    Assume that the image of any pure state density matrix $\Psi$ under $T$ is mixed,
    i.e., $T(\Psi)=\lambda\rho_1+(1-\lambda)\rho_2$ with $\rho_1\neq\rho_2$,
    $\lambda\in(0,1)$. Then by applying $T^{-1}$ to this decomposition we would get a
    nontrivial convex decomposition for $\Psi$ (due to positivity of $T^{-1}$) leading
    to a contradiction. Hence, $T$ and its inverse map pure states onto pure states.
    Furthermore, they are unital, which can again be seen by contradiction. So assume
    $T(\one)\neq \one$. Then the smallest eigenvalue of $T(\one)$ satisfies
    $\lambda_{\min}<1$ due to the trace preserving property. If we denote by $\ket{\lambda}$
    a corresponding eigenvector, then $\one - \frac{\lambda_{\min}+1}{2}\,T^{-1}(\ket{\lambda}\!\bra{\lambda})$
    is a positive operator, but its image under $T$ would no longer be positive.
    Therefore we must have $T(\one)=\one$.

    Every unital positive trace preserving map is contractive with respect to the
    Hilbert--Schmidt norm. As this holds for both $T$ and $T^{-1}$ we have that
    $\forall A\in\mcM_d:\ \|T(A)\|_2=\|A\|_2$, i.e., $T$ acts unitarily on the
    Hilbert--Schmidt Hilbert space. In particular, it preserves the Hilbert--Schmidt
    scalar product $\tr[T(A)T(B)^\dagger]=\tr[AB^\dagger]$. Applying this to pure states
    $A=\ket{\phi}\!\bra{\phi}$ and $B=\ket{\psi}\!\bra{\psi}$ shows that $T$ gives rise
    to a mapping of the Hilbert space onto itself which preserves the value of
    $|\braket{\phi}{\psi}|$. By Wigner's theorem (Thm.~1.1) this has to be either
    unitary or anti-unitary. If $T$ is a unitary conjugation then
    $\det T=\det(U\otimes \overline{U})=1$ due to Eq.~(2.21). Since every anti-unitary
    is unitarily equivalent to complex conjugation, we get that $T$ is in this case a
    matrix transposition $T(A)=A^{\mathsf T}$ (up to unitary equivalence). The determinant
    of the matrix transposition is easily seen in the Gell--Mann basis of $\mcM_d$
    (see Hermitian operator bases in Sec.~2.3). That is, we take basis elements $F_\alpha$
    of the form $\sigma_x/\sqrt{2}$, $\sigma_y/\sqrt{2}$ for $\alpha=1,\dots,d^2-d$ and
    diagonal for $\alpha=d^2-d+1,\dots,d^2$. In this basis matrix transposition is
    diagonal and has eigenvalues $1$ and $-1$ where the latter appears with multiplicity
    $d(d-1)/2$. This means that matrix transposition has determinant minus one iff
    $d(d-1)/2$ is odd, which is equivalent to $\lfloor d/2\rfloor$ being odd.
\end{proof}

\begin{corollary}[Monotonicity of the determinant]
    Let $T,T'$ be two positive and trace-preserving linear maps on $\mcM_d$. Then the
    determinant is decreasing in magnitude under composition, i.e.,
    $|\det T|\ge |\det(TT')|$ where equality holds iff $T'$ is a unitary conjugation,
    a matrix transposition or $\det T=0$.
\end{corollary}

\begin{corollary}[Positive invertible maps]
    Let $T:\mcM_d\to\mcM_d$ be a trace-preserving and positive linear map. Then the
    inverse $T^{-1}$ is positive, too iff $T$ is a unitary conjugation or matrix
    transposition.
\end{corollary}

Note that the inverse (if it exists) of a trace-preserving linear map is
automatically trace-preserving. Hence, the above statement implies that the
inverse of a quantum channel is again a quantum channel iff it describes a
unitary time evolution.

One might wonder whether completely positive maps can have negative determinants.
The following simple example answers this question yes. It is built
up on the map $\rho\mapsto \rho^{T_c}$ which transposes the corners of $\rho\in\mcM_d$,
i.e., $(\rho^{T_c})_{k,l}$ is $\rho_{l,k}$ for the entries $(k,l)=(1,d),(d,1)$ and
remains $\rho_{k,l}$ otherwise. Note that for $d=2$ this is the ordinary matrix
transposition.

\begin{example}[Negative determinant]
    The map $T:\mcM_d\to\mcM_d$ defined by
    \begin{equation}\tag{6.23}\label{eq:6.23}
        T(\rho)=\frac{\rho^{T_c} + \one\,\tr[\rho]}{1+d}
    \end{equation}
    is trace preserving, completely positive with Kraus rank $d^2-1$ and has determinant
    \[
        \det T = -(d+1)^{1-d^2}.
    \]
    For $d=2$ the channel is entanglement breaking and can be written as
    \begin{equation}\tag{6.24}\label{eq:6.24}
        T(\rho)=\frac{1}{3}\sum_{j=1}^6
        \ket{\overline{\xi_j}}\!\bra{\xi_j} \rho\,\ket{\xi_j}\!\bra{\overline{\xi_j}},
    \end{equation}
    where the six $\xi_j$ are the normalized eigenvectors of the three Pauli matrices.
    Moreover, the example given in Eq.~\eqref{eq:6.24} minimizes the determinant within
    the set of all quantum channels on $\mcM_2$.
\end{example}

\begin{proof}
    A convenient matrix representation of the channel is given in the generalized
    Gell--Mann basis. Choose $F_1$ as the $\sigma_y/\sqrt{2}$ element corresponding to
    the corners and $F_2=\one/\sqrt{d}$ the only non-traceless element. Then
    $\widehat{T}=\diag[-1,1+d,1,\dots,1]/(d+1)$ leading to $\det T=-(d+1)^{1-d^2}$.

    For complete positivity we have to check positivity of the Jamiolkowski state $\tau$.
    The corner transposition applied to a maximally entangled state leads to one negative
    eigenvalue $-1/d$. This is, however, exactly compensated by the second part of the map
    such that $\tau\ge 0$ with rank $d^2-1$.

    The representation for $d=2$ is obtained from $\tr[A\,T(B)] = \tr[(A\otimes B^{\mathsf T})\tau]/d$
    by noting that in this case $\tau$ is proportional to the projector onto the symmetric
    subspace which in turn can be written as
    $\frac{1}{2}\sum_j \ket{\xi_j}\!\bra{\xi_j}^{\otimes 2}$ in agreement with the given Kraus
    representation of the channel.

    In order to see that the determinant is minimized for $d=2$, recall from Sec.~2.4
    that any trace-preserving $T$ can be conveniently represented in terms of the real
    $4\times 4$ matrix $\widehat{T}_{ij}:=\tr[\sigma_i T(\sigma_j)]/2$ where the $\sigma_i$
    are identity and Pauli matrices:
    \begin{equation}\tag{6.25}\label{eq:6.25}
        \widehat{T}=
        \begin{pmatrix}
            1 &0\\
            v &\Lambda
        \end{pmatrix},
    \end{equation}
    where $v\in\RR^3$, $\Lambda$ is a $3\times 3$ matrix and $\det\Lambda=\det T$. To simplify
    matters we can diagonalize $\Lambda$ by special orthogonal matrices
    $O_1\Lambda O_2=\diag\{\lambda_1,\lambda_2,\lambda_3\}$ corresponding to unitary operations
    before and after the channel (see Exp.~1.1 for the relation between $\mathrm{SO}(3)$ and
    $\mathrm{SU}(2)$). Obviously, this does neither change the determinant, nor complete
    positivity. For the latter it is necessary that $\vec\lambda$ is contained in a tetrahedron
    spanned by the four corners of the unit cube with $\lambda_1\lambda_2\lambda_3=1$.
    Fortunately, all these points can indeed be reached by unital channels ($v=0$) for which
    this criterion becomes also sufficient for complete positivity. By symmetry we can restrict
    our attention to one octant and reduce the problem to maximizing $p_1p_2p_3$ over all
    probability vectors $\vec p$ yielding $\vec p=(\tfrac13,\tfrac13,\tfrac13)=(\lambda_1,-\lambda_2,\lambda_3)$.
    Hence, the minimal determinant is $-(\tfrac13)^3$ and the corresponding channel can easily
    be constructed from $\widehat{T}$ as $T:\rho\mapsto \frac13(\rho^{\mathsf T}+\one)$.
\end{proof}

\begin{proposition}[Positive determinant for small Kraus rank]
    If $T:\mcM_d\to \mcM_d$ is a completely positive linear map with Kraus rank at most
    two, then $\det T\ge 0$.
\end{proposition}

\begin{proof}
    If $A,B\in\mcM_d$ are Kraus operators then
    $\widehat{T}=A\otimes \overline{A}+B\otimes \overline{B}$ is a matrix representation
    of $T$ (see Eq.~(2.21)). Assume for the moment that $\det A\neq 0$. Then due to
    multiplicativity of the determinant
    \begin{equation}\tag{6.26}\label{eq:6.26}
        \det T
        = \det(A\otimes \overline{A})
        \det\!(\one + A^{-1}B\otimes \overline{A^{-1}B})
        \ge 0.
    \end{equation}
    In order to understand the inequality note that if $\{\lambda_j\}$ are the eigenvalues
    of $A^{-1}B$, then $\{1+\lambda_i\overline{\lambda_j}\}$ are those of
    $\one + A^{-1}B\otimes \overline{A^{-1}B}$. That is, they all come in complex conjugate
    pairs so that their product is positive. The same argumentation holds, of course, for the
    $A\otimes \overline{A}$ term. The proof is completed by using that the set of maps with
    $\det A\neq 0$ is dense and that $\det T$ is continuous.
\end{proof}

\begin{proposition}[Determinant and Choi--Jamiolkowski operator]
    Let $T:\mcM_d\to \mcM_d$ be a linear map and $\tau$ the corresponding Choi--Jamiolkowski
    operator. Then
    \begin{equation}\tag{6.27}\label{eq:6.27}
        |\det T| \le \tr[\tau^\dagger \tau]^{d^2/2}.
    \end{equation}
\end{proposition}

\begin{proof}
    To relate the `purity' to the determinant we exploit that by Eq.~(2.22)
    $\widehat{T}=d\,\tau^\Gamma$ is a matrix representation of the map $T$, so that
    \begin{equation}\tag{6.28}\label{eq:6.28}
        \tr[\tau^\dagger \tau]
        = \frac{1}{d^2}\tr[\widehat{T}^\dagger \widehat{T}]
        = \frac{1}{d^2}\sum_{i=1}^{d^2} s_i^2,
    \end{equation}
    where the $s_i$ are the singular values of $\widehat{T}$ (and thus $T$). From this
    Eq.~\eqref{eq:6.27} is obtained via the geometric--arithmetic mean inequality together
    with the fact that $|\det T|=\prod_i s_i$.
\end{proof}

% ------------------------------------------------------------
\subsection{Irreducible maps and Perron--Frobenius theory}
% ------------------------------------------------------------

Perron--Frobenius theory, roughly speaking, deals with the dominant eigenvalues
and eigenvectors of (element-wise) non-negative matrices (see Exp.\ \dots), i.e.,
in particular with those of classical channels. In order to make the theory work
the considered maps have to fulfill a generic condition like irreducibility,
primitivity or strict positivity (in order of increasing restriction).

In this section we will discuss the generalization of Perron--Frobenius theory to
the context of positive maps on $\mcM_d(\CC)$. The following theorem defines
irreducible positive maps and shows that this can be done in various equivalent
ways:

\begin{theorem}[Irreducible positive maps]
    Let $T:\mcM_d(\CC)\to\mcM_d(\CC)$ be a positive linear map. The following properties
    are equivalent:
    \begin{enumerate}
        \item \textbf{Irreducibility:} If $P\in\mcM_d(\CC)$ is a Hermitian projector such
              that $T(P\mcM_dP)\subseteq P\mcM_dP$, then $P\in\{0,\one\}$.
        \item For every non-zero $A\ge 0$ we have $(\id + T)^{d-1}(A)>0$.
        \item For every non-zero $A\ge 0$ and every strictly positive $t\in\RR$ we have
              $\exp[tT](A)>0$.
        \item For every orthogonal pair of non-zero, positive semi-definite $A,B\in\mcM_d(\CC)$
              there is an integer $t\in\{1,\dots,d-1\}$ such that $\tr[B\,T^t(A)]>0$.
    \end{enumerate}
\end{theorem}

\begin{proof}
    \begin{itemize}
        \item[1.\ $\Rightarrow$ 2.]
              From $T(A)\ge 0$ we get an inclusion for the kernels
              $\ker(\id+T)(A)\subseteq \ker A$. Suppose equality holds in this inclusion, then
              $\supp(\id+T)(A)\subseteq \supp A$. Therefore $T(P\mcM_dP)\subseteq P\mcM_dP$ if
              we take $P$ the Hermitian projector onto the support space of $A$. By 1.
              (irreducibility) this can only be if $A>0$. Application of $(\id+T)$ thus has to
              increase the rank until there is no kernel left, which happens after $d-1$ steps
              at the latest.

        \item[2.\ $\Rightarrow$ 3.]
              Looking at the series expansion of both $(\id+T)^{d-1}(A)$ and $\exp[tT](A)$ in
              terms of powers of $T$ we realize that (i) all terms are positive, and (ii) all
              terms appearing in the first expansion are also present in the second one.
              Therefore we can bound $\exp[tT](A)\ge c(\id+T)^{d-1}(A)$ with some non-zero
              positive constant $c$.

        \item[3.\ $\Rightarrow$ 1.]
              Suppose $T$ is reducible, i.e., there is a proper projection $P\neq 0,\one$ for
              which $T(P)\le cP$ for some scalar $c>0$. Then $\exp[tT](P)\le \exp[tc]P$, so $T$
              does not fulfill 3.

        \item[4.\ $\Rightarrow$ 1.]
              If $T$ is reducible by a proper projection $P$, then $\tr[(\one-P)T^t(P)]=0$ for
              all $t\in\NN$.

        \item[2.\ $\Rightarrow$ 4.]
              Expanding $\tr[B(\id+T)^{d-1}(A)]>0$ and using $\tr[BA]=0$ together with positivity
              of all terms in the expansion, we get that $\tr[B\,T^t(A)]>0$ for at least one
              $t\le d-1$.
    \end{itemize}
\end{proof}

Note that $T$ is irreducible iff the dual map $T^*$ is irreducible. If a proper
projection $P$ `reduces' $T$ in the sense that $\tr[(\one-P)T(P)]=0$, then $\one-P$
reduces $T^*$.

In order to relate irreducibility to spectral properties of a positive map $T$ it
is useful to consider the following functionals defined on the cone of positive
semi-definite operators:
\begin{equation}\tag{6.29}\label{eq:6.29}
    r(X) := \sup\{\lambda\in\RR \mid (T-\lambda\id)(X)\ge 0\},
\end{equation}
\begin{equation}\tag{6.30}\label{eq:6.30}
    \tilde r(X) := \inf\{\lambda\in\RR \mid (T-\lambda\id)(X)\le 0\}.
\end{equation}
We are especially interested in the maxima
$r := \sup_{X\ge 0} r(X)$ and $\tilde r := \sup_{X\ge 0}\tilde r(X)$ which
obviously satisfy $r\ge \tilde r$ and actually coincide for irreducible maps:

\begin{theorem}[Spectral radius of irreducible maps]
    Let $T:\mcM_d(\CC)\to \mcM_d(\CC)$ be an irreducible positive map. Then
    \begin{enumerate}
        \item For the quantities introduced below Eqs.~\eqref{eq:6.29},\eqref{eq:6.30}
              we have $r=\tilde r$.
        \item $r$ is a non-degenerate eigenvalue of $T$ and the corresponding eigenvector
              is strictly positive, i.e., $T(X)=rX>0$.
        \item If there is any $\lambda>0$ which is an eigenvalue of $T$ with positive
              eigenvector, i.e., $T(Y)=\lambda Y\ge 0$, then $\lambda=r$.
        \item $r$ is the spectral radius of $T$.
    \end{enumerate}
\end{theorem}

\begin{proof}
    First note that all suprema/infima in Eqs.~\eqref{eq:6.29},\eqref{eq:6.30} and for
    defining $r,\tilde r$ are attained since the respective sets are compact or can be
    made so, e.g., by imposing $\tr[X]=1$.

    We begin with showing that $r$ is attained for a strictly positive $X$ and that
    $T(X)=rX$. To this end consider any non-zero $X\ge 0$ for which $\lambda:=r(X)>0$.
    The identity
    \begin{equation}\tag{6.31}\label{eq:6.31}
        (T+\id)^{d-1}(T-\lambda\id)(X) = (T-\lambda\id)(T+\id)^{d-1}(X),
    \end{equation}
    shows two things. First, the supremum $\sup_{X\ge 0} r(X)$ is attained for a strictly
    positive $X$, since by Thm.~6.2 $(T+\id)^{d-1}(X)>0$. Second, the $X$ achieving
    $r(X)=r$ must satisfy $(T-r\id)(X)=0$, since otherwise the expression in
    Eq.~\eqref{eq:6.31} would be positive definite so that a larger multiple of the identity
    map could be subtracted---contradicting the supposed maximality. As $r(X)=\tilde r(X)$
    for every eigenvector $X\ge 0$, the just mentioned observation together with the
    inequality $r\ge \tilde r$ proves item~1.

    To prove 2.\ we still need to show non-degeneracy. Suppose there would be an
    eigenvector $X'$ corresponding to $r$ which is not a multiple of $X$. By taking the
    Hermitian (or $i$ times the anti-Hermitian) part we can assume $X'=X'^\dagger$.
    As $X>0$ there is a non-zero $c\in\RR$ such that $0\le X+cX'$ has a kernel. This is,
    however, in conflict with
    \begin{equation}\tag{6.32}\label{eq:6.32}
        0 < (T+\id)^{d-1}(X+cX') = (r+1)^{d-1}(X+cX'),
    \end{equation}
    so that $X'$ must be a multiple of $X$. Hence, $r$ is indeed non-degenerate.

    For the proof of 3.\ assume that $Y\ge 0$ is an eigenvector of $T$ corresponding to
    an eigenvalue $\lambda>0$. Using that $r$ is an eigenvalue of the dual map $T^*$ with a
    strictly positive eigenvector $X_0>0$ we get
    \begin{equation}\tag{6.33}\label{eq:6.33}
        r\,\tr[X_0Y] = \tr[T^*(X_0)Y] = \tr[X_0T(Y)] = \lambda\,\tr[X_0Y].
    \end{equation}
    Since $\tr[X_0Y]>0$ this implies $r=\lambda$.

    Finally, in order to prove 4.\ we define a positive map
    \[
        T'(\,\cdot\,)
        := \frac{1}{r}\,X^{-1/2}\,T\!(X^{1/2}(\,\cdot\,)X^{1/2})\,X^{-1/2}
    \]
    with $X>0$ being the positive eigenvector corresponding to the eigenvalue $r$ of $T$.
    Note that up to the factor $1/r$ the maps $T$ and $T'$ differ only by a similarity
    transformation, i.e., their spectra coincide up to that factor. In particular the
    spectral radii satisfy $\varrho(T')=\varrho(T)/r$. Since $T'$ is a unital positive map
    we have $\varrho(T')=1$ by Prop.~6.1 and thus $\varrho(T)=r$.
\end{proof}

\begin{proposition}[Similarity transformations preserving irreducibility]
    Let $T:\mcM_d(\CC)\to\mcM_d(\CC)$ be an irreducible positive map. Then for arbitrary
    constants $c>0$ and invertible matrices $C\in\mcM_d(\CC)$ the maps
    \begin{equation}\tag{6.34}\label{eq:6.34}
        T'(\,\cdot\,) := c\,C^{-1}T(C\,\cdot\, C^\dagger)C^{-\dagger},
    \end{equation}
    are irreducible as well.
\end{proposition}

\begin{proof}
    Assume that a projection $Q$ would `reduce' $T'$. Then the projection $P$ onto the
    support of $CQC^\dagger$ would reduce $T$. By invertibility of $C$ we have that $P$
    and $Q$ have the same rank and thus $T'$ can only be reducible if $T$ is.
\end{proof}

\begin{theorem}[Irreducibility from spectral properties]
    Let $T:\mcM_d(\CC)\to\mcM_d(\CC)$ be a positive map with spectral radius $\varrho$.
    The following are equivalent:
    \begin{enumerate}
        \item $T$ is irreducible.
        \item The spectral radius $\varrho$ is a non-degenerate eigenvalue and the corresponding
              right and left eigenvectors are positive definite (i.e., $T(X)=\varrho X>0$ and
              $T^*(Y)=\varrho Y>0$).
    \end{enumerate}
\end{theorem}

\begin{proof}
    $1.\Rightarrow 2$.\ is given by Thm.~6.3 when applied to $T$ and $T^*$.

    For the converse $2.\Rightarrow 1$.\ we apply Eq.~\eqref{eq:6.34} with $c=1/\varrho$
    and $C=Y^{-1/2}$. In this way $T'$ becomes trace preserving with $T'(X')=X'>0$,
    $X':=\sqrt{Y}\,X\sqrt{Y}$. Moreover, the non-degeneracy translates to that of the
    eigenvalue $1$ of $T'$. Now assume that $T'$ (and thus $T$) would be reducible, i.e.,
    there is a Hermitian projection $P\notin\{0,\one\}$ such that $T':A\to A$ with
    $A:=P\mcM_d(\CC)P$. In this case $T'$ would have a fixed point density operator
    $\rho\in A$ for which $T'(\rho)=\rho\ge 0$ (see Sec.~6.4). Hence, the eigenvalue $1$
    of $T'$ would have to be degenerate since $\rho$ is not a multiple of $X'$ because
    $X'\notin A$ (as $X'>0$).
\end{proof}

\begin{corollary}[Time-average and ergodicity]
    Let $T:\mcM_d(\CC)\to\mcM_d(\CC)$ be a positive, trace-preserving linear map. Then
    the following are equivalent:
    \begin{enumerate}
        \item $T$ is irreducible,
        \item There is a unique state with density matrix $\sigma>0$ such that for every
              density matrix $\rho\in\mcM_d(\CC)$ we have
              \begin{equation}\tag{6.35}\label{eq:6.35}
                  \lim_{N\to\infty}\frac{1}{N}\sum_{t=1}^N T^t(\rho) = \sigma.
              \end{equation}
    \end{enumerate}
\end{corollary}

\begin{proof}
    Note that Eq.~\eqref{eq:6.35} is nothing but $T_\infty(\rho)=\sigma$, where $T_\infty$ is
    the projection onto the eigenspace of $T$ with eigenvalue one as given in
    Eq.~\eqref{eq:6.14}. The equivalence $1.\Leftrightarrow 2$.\ thus follows from the
    equivalence in Thm.~6.4.
\end{proof}

\begin{theorem}[Spectral radius and positive eigenvectors]
    Let $T:\mcM_d(\CC)\to\mcM_d(\CC)$ be a positive map with spectral radius $\varrho$.
    Then $\varrho$ is an eigenvalue and there is a positive semi-definite $X\in\mcM_d(\CC)$
    such that $T(X)=\varrho X$.
\end{theorem}

The peripheral spectrum is the set of eigenvalues whose magnitude is equal to the
spectral radius, i.e., equal to one if we are considering unital positive maps
(Prop.~6.1). As we saw in Prop.~6.2 the properties positivity and unitality together
imply that all eigenvalues in the peripheral spectrum have trivial (i.e.,
one-dimensional) Jordan blocks. If we add irreducibility and the assumed validity
of the Schwarz inequality in Eq.~(5.2) we can say much more:

\begin{theorem}[Peripheral spectrum of irreducible Schwarz maps]
    Let $T:\mcM_d(\CC)\to\mcM_d(\CC)$ be an irreducible positive unital map which fulfills
    the Schwarz inequality in Eq.~(5.2). Denote by $S:=\spec(T)\cap \exp(i\RR)$ the
    peripheral spectrum of $T$. Then
    \begin{enumerate}
        \item There is an integer $m\in\{1,\dots,d^2\}$ such that
              $S=\{\exp(2\pi i k/m)\}_{k\in\ZZ_m}$.
        \item All eigenvalues in $S$ are non-degenerate.
        \item There is a unitary $U$ such that for all $k\in\NN$: $T(U^k)=\gamma^k U^k$
              where $\gamma:=\exp(2\pi i/m)$.
        \item $U$ has spectral decomposition $U=\sum_{k\in\ZZ_m}\gamma^k P_k$ where the
              spectral projections $P_k$ satisfy $T(P_{k+1})=P_k$.
    \end{enumerate}
\end{theorem}

\begin{proof}
    Take any eigenvalue $\lambda\in S$ with corresponding eigenvector $U\in\mcM_d(\CC)$,
    i.e., $T(U)=\lambda U$. Then the Schwarz inequality reads $U^\dagger U\le T(U^\dagger U)$.
    Since $T$ is irreducible by assumption we can apply the tools in the proof of Thm.~6.3,
    in particular that $r(U^\dagger U)=1$, which implies that equality has to hold in the
    Schwarz inequality and that $U$ can be rescaled to be a unitary. Moreover, by Thm.~5.3 we have
    \begin{equation}\tag{6.36}\label{eq:6.36}
        \overline{\lambda}\,U^\dagger T(X)=T(U^\dagger X),
        \qquad \forall X\in\mcM_d(\CC).
    \end{equation}
    Applying this to the case where $X$ is another eigenvector of $T$ w.r.t.\ an eigenvalue
    in $S$ shows that $S$ is closed under multiplication and thus a group. Since there are at
    most $d^2$ elements in $S$ we must have $S=\exp(2\pi i \ZZ_m/m)$ for some integer $m\le d^2$.

    Now assume that $U$ and $V$ are both eigenvectors w.r.t.\ the same eigenvalue in $S$.
    Then $T(U^\dagger V)=U^\dagger V$ and, since the eigenvalue one of an irreducible unital
    positive map is non-degenerate (Prop.~6.1 and Thm.~6.3), we must have $U\propto V$.
    Hence, all eigenvalues in $S$ are simple.

    If $U$ is the unitary eigenvector corresponding to the eigenvalue $\gamma:=\exp(2\pi i/m)$,
    then the group property gives $T(U^k)=\gamma^k U^k$ for all $k\in\NN$. In particular
    $T(U^m)=U^m$ which has to be equal to $\one$ by non-degeneracy of the eigenvalue one.
    Therefore we can write $U=\sum_{k\in\ZZ_m}\gamma^k P_k$ with orthogonal spectral projections
    $P_k$. This gives
    \begin{equation}\tag{6.37}\label{eq:6.37}
        \sum_{k\in\ZZ_m}\gamma^{k+1}P_k = \gamma U = T(U) = \sum_{k\in\ZZ_m}\gamma^k T(P_k).
    \end{equation}
    Due to $T(U^k)=\gamma^kU^k$, $T$ acts as an automorphism on the $*$-algebra generated by
    $U$. Projections are thus mapped onto projections and by comparing the left and right hand
    side in Eq.~\eqref{eq:6.37} we get $T(P_{k+1})=P_k$ from the uniqueness of the spectral decomposition.
\end{proof}

\begin{proposition}[Covariance and eigensystem of irreducible maps]
    Let $T:\mcM_d(\CC)\to\mcM_d(\CC)$ be an irreducible positive and unital linear map which
    fulfills the Schwarz inequality. Following Thm.~6.6 denote by $\{\gamma^k\}_{k=1,\dots,m}$
    the peripheral spectrum of $T$ and by $U=\overline{\gamma}\,T(U)$ the respective unitary
    eigenvector.
    \begin{enumerate}
        \item For every $X\in\mcM_d(\CC)$ we have $T(U^\dagger XU)=U^\dagger T(X)U$.
        \item If $X$ is an eigenvector with $T(X)=\mu X$, then
              \begin{equation}\tag{6.38}\label{eq:6.38}
                  T(U^\dagger X)=\overline{\gamma}\mu\,U^\dagger X,
                  \qquad
                  T(XU)=\gamma\mu\,XU.
              \end{equation}
    \end{enumerate}
\end{proposition}

\begin{proof}
    Both statements follow directly from Eq.~\eqref{eq:6.36} and its adjoint.
\end{proof}


% === Content below merged from r0 (r1 was missing proofs in later sections) ===
% ------------------------------------------------------------
\subsection{Primitive maps}
% ------------------------------------------------------------

We will continue the discussion on Perron--Frobenius theory and direct the focus
to primitive maps. A primitive map is an irreducible positive map with trivial
peripheral spectrum. Similar to the irreducible case primitive maps can,
however, be characterized by a number of equivalent properties each of which
may be used as a definition:

\begin{theorem}[Primitive maps]
    Let $T:\mcM_d(\CC)\to\mcM_d(\CC)$ be a positive and trace-preserving linear map.
    Then the following statements are equivalent:
    \begin{enumerate}
        \item There exists an $n\in\NN$ such that for every density matrix
              $\rho\in\mcM_d(\CC)$ we have $T^n(\rho)>0$.
        \item $T^k$ is irreducible for every $k\in\NN$ and the limit $T_\phi$
              (see Prop.~6.3) is irreducible as well.
        \item For all density matrices $\rho$ the limit $\lim_{k\to\infty} T^k(\rho)$
              exists, is independent of $\rho$ and given by a positive definite density
              matrix $\rho_\infty$.
        \item $T$ has trivial peripheral spectrum (i.e., counting algebraic
              multiplicities, there is only one eigenvalue of magnitude one) and the
              corresponding eigenvector is positive definite.
    \end{enumerate}
\end{theorem}

\begin{proof}
    $1.\Rightarrow 2$. can be proven by contradiction: assume for instance
    $T_\phi=\lim_{i\to\infty}T^{n_i}$ is reducible. Then there is a density matrix
    $\rho_0$ which has a kernel and for which
    $\rho_0=T_\phi(\rho_0)=\lim_i T^n\!(T^{n_i-n}(\rho))$ which would imply
    that $T^n(\rho)\not>0$ for some $\rho$. Similarly, if $T^k$ and therefore also
    $T^{km}$ is reducible for all $m\in\NN$, then there cannot be any $n$ for which
    1.\ holds.

    $2.\Rightarrow 4$.: If $T_\phi$ has a non-degenerate eigenvalue one, then, by
    definition of $T_\phi$, the map $T$ has trivial peripheral spectrum. Moreover,
    since the fixed point set of $T$ is included in the one of $T_\phi$ which in
    turn contains only a single positive definite element, this fixed point has to
    be the one of $T$ as well.

    $4.\Rightarrow 3$.: Since $T$ has only one eigenvalue of magnitude one with
    corresponding positive definite eigenvector $\rho_\infty=T(\rho_\infty)$ we have
    $\lim_{k\to\infty}T^k=T_\phi$ and $T_\phi(\rho)=\rho_\infty \tr[\rho]$ is a
    projection.

    $3.\Rightarrow 1$.: First note that 3.\ implies that the eigenvalue of $T$ which
    is second largest in magnitude, denote it by $\lambda_2$, satisfies
    $|\lambda_2|<1$. Assume that $T^n(\rho)$ would have a kernel with eigenvector
    $\psi$. Then
    \begin{align}
        \lambda_{\min}(\rho_\infty)
         &\le \bigl|\langle\psi|T^n(\rho)-\rho_\infty|\psi\rangle\bigr|
        \le \lVert T^n(\rho)-\rho_\infty\rVert, \tag{6.39}\\
         &= \lVert (T^n-T_\infty)(\rho-\rho_\infty)\rVert
        \le \mu^n c\,\lVert\rho-\rho_\infty\rVert, \tag{6.40}
    \end{align}
    where $\lambda_{\min}$ denotes the smallest eigenvalue, $\mu$ is such that
    $1>\mu>|\lambda_2|$ and $\lVert\cdot\rVert$ is the operator norm. The constant
    $c>0$ depends on $T$ but is independent of $n$ (see the discussion around
    Thm.~8.23 for more details). Thus if $n$ is sufficiently large, then
    Eqs.~\eqref{6.39},\eqref{6.40} would lead to a contradiction for every $\rho$.
    Hence, there is a finite $n$ such that $T^n(\rho)>0$ for all $\rho$.
\end{proof}

If we interpret $T^t(\rho)$ for $t\in\NN$ as discrete time-evolution, then
primitive maps are exactly those for which the evolution eventually converges to
a stationary state which is strictly positive (i.e., it may in some contexts be
interpreted as a finite temperature Gibbs state) and independent of the initial
state.

Note that if $T$ in addition satisfies the Schwarz inequality as required by
Thm.~6.6, then we can replace point~2.\ in Thm.~6.7 by the statement that $T^k$
is irreducible for all $k<d^2$.

Moreover, if a map is completely positive, we can give three more equivalent
characterizations of primitivity:

\begin{theorem}[Completely positive primitive maps]
    Let $T:\mcM_d(\CC)\to\mcM_d(\CC)$ be a completely positive and trace-preserving
    linear map with Kraus decomposition $T(\cdot)=\sum_i K_i\,\cdot\,K_i^\dagger$.
    Denote by
    \[
        K_m := \spn\left\{\prod_{k=1}^m K_{i_k}\right\}
    \]
    the complex linear span of all degree-$m$ monomials of Kraus operators. Then the
    following are equivalent:
    \begin{enumerate}
        \item $T$ is primitive.
        \item There exists an $n\in\NN$ such that for every non-zero $\psi\in\CC^d$ and
              all $m\ge n$: $K_m\ket{\psi}=\CC^d$.
        \item There exists a $q\in\NN$ such that for all $m\ge q$: $K_m=\mcM_d(\CC)$.
        \item There exists a $q\in\NN$ such that for all $m\ge q$:
              $(T^m\otimes \id_d)(\ket{\Omega}\!\bra{\Omega})>0$, where
              $\Omega\in\CC^d\otimes \CC^d$ denotes a maximally entangled state.
    \end{enumerate}
    The numbers $q$ appearing in 3., 4.\ can be chosen equal and if $q$ and $n$ are
    chosen minimal, then $n\le q$.
\end{theorem}

\begin{proof}
    First consider item~2.\ and observe that it obviously holds for all $m\ge n$ if
    it holds for some $n$ since we can always incorporate $m-n$ Kraus operators in
    $\psi$. Similar holds for item~4. In order to see this, define
    $\tau_q:=(T^q\otimes \id_d)(\ket{\Omega}\!\bra{\Omega})$ and assume $\tau_q>0$.
    Then there is in particular a constant $c>0$ such that $\tau_q\ge c\tau_{q-1}$
    and therefore
    \begin{equation}\tag{6.41}
        \tau_{q+1}
        = (T\otimes \id)(\tau_q)
        \ge c(T\otimes \id)(\tau_{q-1})
        = c\tau_q>0.
    \end{equation}
    By induction the same holds then true for all integers larger than $q$.

    \noindent
    $1.\Leftrightarrow 2$.:
    Let us introduce $K_I:=\prod_{k=1}^n K_{i_k}$ with multi-index
    $I=(i_1,\dots,i_n)$. If $T$ is primitive, then by Thm.~6.7 there exists an $n$
    such that $\sum_I K_I\ket{\psi}\!\bra{\psi}K_I^\dagger$ has full rank for every
    $\psi$. Thus, for every $\phi$ there is an $I$ so that
    $\langle\phi|K_I|\psi\rangle\neq 0$ and therefore
    $\spn\{K_I\ket{\psi}\}=\CC^d$. Conversely, if the latter is true, then for every
    $\phi$ there will be a unitary $U$ so that
    $\ket{\phi}\propto \sum_I U_{JI}K_I\ket{\psi}=:K'_J\ket{\psi}$. Since $\{K_I\}$
    and $\{K'_J\}$ both provide a Kraus representation of $T^n$ (see Thm.~2.1) the
    full rank property is guaranteed. Following Thm.~6.7 $T$ then has to be
    primitive.

    \noindent
    $3.\Leftrightarrow 4$.:
    Following the above lines and denoting by $K_I$ the Kraus operators of $T^q$ it
    is clear that $(T^q\otimes \id_d)(\ket{\Omega}\!\bra{\Omega})>0$ is equivalent to
    $\spn\{(K_I\otimes \one_d)\ket{\Omega}\}=\CC^d\otimes \CC^d$. From here the
    equivalence between item~3.\ and 4.\ follows due to the fact that the relation
    $K\leftrightarrow (K\otimes \one)\ket{\Omega}$ is a linear one-to-one mapping
    between $\mcM_d(\CC)$ and $\CC^d\otimes \CC^d$.

    \noindent
    $4.\Rightarrow 1$.:
    If the Choi--Jamiolkowski matrix
    $\tau:=(T^q\otimes \id_d)(\ket{\Omega}\!\bra{\Omega})$ has full rank, then so
    does
    \begin{equation}\tag{6.42}
        T^q(\ket{\psi}\!\bra{\psi})
        = d(\one_d\otimes \bra{\overline{\psi}}) \tau\,(\one_d\otimes \ket{\overline{\psi}})
        \qquad \forall \psi\in\CC^d.
    \end{equation}
    This also shows that $n=q$ is sufficient. So if we choose both $n$ and $q$
    minimal, then $n\le q$ in general.

    \noindent
    $1.\Rightarrow 4$.:
    The proof for this implication is entirely parallel with
    Eqs.~\eqref{6.39},\eqref{6.40} where $q$ now plays the role of $n$: just replace
    $\rho$ by $\Omega$, $\rho_\infty$ by $\rho_\infty\otimes \one/d$ and $T$ by
    $T\otimes \id$. Also the constant $c$ now depends on $T\otimes \id$, but the
    argumentation remains the same.
\end{proof}

Note that all the above relations for both positive and completely positive maps
hold as well for maps which are not trace-preserving\footnote{The limit
    $\lim_{k\to\infty}T^k$ has then to be normalized by the $k$th power of the
    spectral radius of $T$.} so that the notion of primitivity is well-defined for
arbitrary positive maps.

Finally, we will show that the number $q$ (and therefore also $n$) appearing in
Thm.~6.8 can be bounded by the dimension $d$ of the system. These bounds can be
seen as a quantum counterpart to Wielandt's inequality, which does the same for
the classical case of primitive stochastic matrices. We need two preparatory
Lemmas before the main result. As before $K_n$ denotes the complex linear span
of the Kraus operators of $T^n$. We will frequently use that every $K\in K_n$ is,
up to a scalar multiple, a possible Kraus operator of $T^n$. This follows from
the freedom in the Kraus decomposition (Thm.~2.1, point~4.).

\begin{lemma}
    Let $T$ be a primitive quantum channel on $\mcM_d(\CC)$ with Kraus rank $k$. Then,
    there is a natural number $n\le d^2-k+1$ and a $K\in K_n$ such that $\tr[K]\neq 0$.
\end{lemma}

\begin{proof}
    Let us denote by $R_n$ the complex linear span of all $K_m$ with $m\le n$. We have
    to show that: (*) for any $n\in\NN$, if $\dim(R_n)<d^2$, then
    $\dim(R_{n+1})>\dim(R_n)$. Since $\dim(R_1)=k$, by iteration we obtain that
    $R_{d^2-k+1}=\mcM_d(\CC)$. This implies that a linear combination of the elements
    of $K_n$ with various $n\le d^2-k+1$ must be equal to the identity, and thus at
    least one of the elements must have non-zero trace. To prove (*) we note that,
    by definition, $R_n\subseteq R_{n+1}$. If they would be equal, then $R_m=R_n$ for
    all $m>n$. Thus, equality can only occur when $\dim(R_n)=d^2$ since otherwise
    the map $T$ would not be primitive.
\end{proof}

\begin{lemma}
    Let $T:\mcM_d(\CC)\to\mcM_d(\CC)$ be a primitive completely positive map such that
    at least one of its Kraus operators, say $K_1$, has a non-zero eigenvalue, i.e.,
    $K_1\ket{\varphi}=\mu\ket{\varphi}$ with $\mu\neq 0$. Then:
    \begin{enumerate}
        \item[(a)] $K_{d-1}\ket{\varphi}=\CC^d$.
        \item[(b)] If $K_1$ is not invertible, then for all $\ket{\psi}\in\CC^d$,
              $\ket{\varphi}\!\bra{\psi}\in K_{d^2-d+1}$.
    \end{enumerate}
\end{lemma}

\begin{proof}
    (a) We define $S_n$ as the span of all $K_m\ket{\varphi}$ with $m\le n$ together
    with $\ket{\varphi}$. If $\dim(S_n)<d$, then $\dim(S_{n+1})>\dim(S_n)$, since
    otherwise the map would not be primitive. Thus, $S_{d-1}=\CC^d$. That is, for all
    $\ket{\phi}\in\CC^d$, there exist matrices $K^{(n)}\in K_{k_n}$, $k_n\le d-1$ such
    that (with $K^{(0)}\propto \one$)
    \begin{equation}\tag{6.43}
        \ket{\phi}
        = \sum_{n=0}^{d-1} K^{(n)}\ket{\varphi}
        = \sum_{n=0}^{d-1} K^{(n)} \frac{K_1^{d-k_n}}{\mu^{d-k_n}}\ket{\varphi},
    \end{equation}
    and thus $\ket{\phi}\in K_{d-1}\ket{\varphi}$.

    (b) We write $K_1$ in the Jordan standard form and divide it into two blocks. The
    first one, of size $\tilde d\times \tilde d$, consists of all Jordan blocks
    corresponding to non-zero eigenvalues, whereas the second one contains all those
    corresponding to zero eigenvalues. We denote by $P$ the (not necessarily
    Hermitian) projector onto the subspace where the first block is supported and by
    $r\le d-\tilde d$ the size of the largest Jordan block corresponding to a zero
    eigenvalue. We have
    \begin{equation}\tag{6.44}
        K_1P=PK_1,
        \qquad
        K_1^r = K_1^r P.
    \end{equation}
    We define $R_n:=PK_n$ and show that $R_{d\tilde d}=P\mcM_d(\CC)$. For all $n\in\NN$,
    $\dim(R_{n+1})\ge \dim(R_n)$. The reason is that for any set of linearly
    independent matrices $K^{(n)}_k\in R_n$, $K_1K^{(n)}_k\in R_{n+1}$ are also linearly
    independent, given that $K_1$ is invertible on its range. By following the
    reasoning of [?, Appendix A] we get that, if $\dim(R_{n+1})=\dim(R_n)=:d'$, then
    $\dim(R_m)=d'$ for all $m>n$, which is incompatible with $T$ being primitive unless
    $d'=\tilde d d$. Thus, for any $\ket{\psi}\in\CC^d$, there exists a $K\in K_{d\tilde d}$
    with
    \[
        \ket{\varphi}\!\bra{\psi}
        = PK
        = \frac{K_1^r PK}{\mu^r}
        = \frac{K_1^r K}{\mu^r}
        \in K_{d\tilde d + r}.
    \]
    By using that $\tilde d\le d-r$ and that $r\ge 1$ (since $K_1$ is supposed to be
    not invertible) we get $\tilde d d + r \le d^2-d+1$, which concludes the proof.
\end{proof}

We have now the necessary tools to prove our main result.

\begin{theorem}[Quantum Wielandt inequality]
    Let $T$ be a primitive quantum channel on $\mcM_d(\CC)$ with Kraus rank $k$. Then
    for the minimal $q$ as it appears in Thm.~6.8 we have that
    \begin{enumerate}
        \item in general $q\le (d^2-k+1)d^2$,
        \item if the span of Kraus operators $K_1$ contains an invertible element, then
              $q\le d^2-k+1$,
        \item if the span of Kraus operators $K_1$ contains an element with at least one
              non-zero eigenvalue, then $q\le d^2$.
    \end{enumerate}
\end{theorem}

\begin{proof}
    \textbf{2.} If there is an invertible element, then it follows from [?, Appendix A,
            Proposition 2] that $\dim K_{n+1}>\dim K_n$ until the full matrix space $\mcM_d(\CC)$ is
    spanned and thus $q\le d^2-k+1$.

    \textbf{1.} Let us denote by $\{K^{(n)}_k\}$ the Kraus operators corresponding to
    $T^n$. According to Lemma~6.2, one of them, say $K^{(n)}_1$, has non-zero trace for
    some $n\le d^2-k+1$ and therefore there exists $\ket{\varphi}$ such that
    $K^{(n)}_1\ket{\varphi}=\mu\ket{\varphi}$ with $\mu\neq 0$. If $K^{(n)}_1$ is
    invertible, then 1.\ is implied by 2., so we can assume that it is not invertible.
    According to Lemma~6.3(b), for all $\ket{\psi}$, $\ket{\varphi}\!\bra{\psi}\in
        K_{(d^2-d+1)n}$; and according to Lemma~6.3(a) $\ket{\chi}\!\bra{\psi}\in K_{nd^2}$.
    This implies that $K_{nd^2}=\mcM_d(\CC)$ and hence the general bound 1.\ follows.

    The argument which proves 3.\ is completely analogous just with $n=1$ as we do,
    by assumption, not have to block in order to get a Kraus operator with non-zero
    eigenvalue.
\end{proof}

In the previous proof we use, for the general case, blocks of $T^n$ for some
$n\le d^2-k+1$ just in order to get any Kraus operator with a non-zero
eigenvalue. This rather clumsy step is not necessary if the dimension is small:

\begin{corollary}
    Let $T$ be a primitive quantum channel on $\mcM_d(\CC)$ with $d=2,3$. Then there
    is always a Kraus operator with non-zero eigenvalue in $K_1$ and thus $q\le d^2$.
\end{corollary}

\begin{proof}
    The fact that $K_1$ has this property for $d=2,3$ stems from the classification
    of nilpotent subspaces~[?]: assume that $K_1$ would be a nilpotent subspace
    within the space of $d\times d$ matrices. Then for $d=2$ its dimension would have
    to be one, so it could not arise from the Kraus operators of a quantum channel.
    Similarly, for $d=3$ there are (up to similarity transformations) two types of
    nilpotent subspaces~[?] of dimension $k>1$: one of dimension $k=3$, the space of
    upper-triangular matrices, whose structure does not allow the trace-preserving
    property, and one of dimension $k=2$ which only leads to quantum channels having
    a (in modulus) degenerate largest eigenvalue. Hence, if $K_1$ is generated by the
    Kraus operators of a primitive quantum channel, then it cannot be nilpotent if
    $d=2,3$.
\end{proof}

For $d\ge 4$ the set of Kraus operators can form a nilpotent subspace even for
primitive quantum channels. However, exploiting the structure of such
subspaces~[?] the general bound in Thm.~6.9 can still be improved.

\setcounter{problem}{11}
\begin{problem}[Optimal Wielandt bound]
Find the smallest integers $n$ and $q$ (as a function of the dimension $d$) for
Thm.~6.8.
\end{problem}

% ------------------------------------------------------------
\subsection{Fixed points}
% ------------------------------------------------------------

\begin{theorem}[Brouwer's fixed point theorem]
    Let $T$ be a continuous map from a non-empty, compact, convex set
    $S\subset \RR^n$ into itself, then there is an $x\in S$ such that $T(x)=x$.
\end{theorem}

As the set of Hermitian matrices in $\mcM_d(\CC)$ is a real vector space and the
density matrices therein form a compact, convex set, Brouwer's fixed point
theorem immediately implies:

\begin{theorem}[Stationary states]
    Every continuous, trace-preserving, positive (not necessarily linear) map
    $T:\mcM_d(\CC)\to\mcM_d(\CC)$ has at least one stationary state. That is, there
    is a density matrix $\rho\in\mcM_d(\CC)$ such that $T(\rho)=\rho$.
\end{theorem}

In the following we are interested in linear maps, which are automatically
continuous.

\begin{proposition}[Positive fixed-points]
    Let $T:\mcM_d(\CC)\to\mcM_d(\CC)$ be a trace-preserving, positive, linear map and
    $X=T(X)$ a fixed point. If $X=\sum_{j=1}^4 c_j P_j$ is the decomposition of $X$
    into four positive operators $P_j\ge 0$ which arises from decomposing $X$ first
    into Hermitian and anti-Hermitian parts and these further into orthogonal
    positive and negative parts, then
    \begin{equation}\tag{6.45}
        T(P_j)=P_j,\qquad j=1,\dots,4,
    \end{equation}
    are fixed points as well.
\end{proposition}

\begin{proof}
    For every Hermiticity preserving linear map we obviously have that the two
    Hermitian operators $X+X^\dagger$ and $i(X-X^\dagger)$ are fixed points if $X$ is
    one. So let us assume that $X$ is Hermitian and that $P_\pm\ge 0$ are its
    orthogonal positive and negative parts, i.e., $X=P_+-P_-$ and $\tr[P_+P_-]=0$.
    The fixed point equation then reads $T(P_+)-T(P_-)=P_+-P_-$ and we want to
    identify $T(P_\pm)$ with $P_\pm$. To this end, denote by $Q$ the projection onto
    the support of $P_+$ and note that the fixed point equation together with the
    trace-preserving property of $T$ implies that
    \begin{align}
        \tr[P_+]
         &= \tr[Q(P_+-P_-)]
        = \tr\!\left[Q\,T(P_+-P_-)\right], \tag{6.46}\\
         &\le \tr[T(P_+)] = \tr[P_+]. \tag{6.47}
    \end{align}
    Since equality has to hold for the inequality in Eq.~\eqref{6.47}, we have to
    have that $T(P_+)$ is supported on $Q$ while $T(P_-)$ is orthogonal to it.
    Consequently, $T$ preserves the orthogonality of $P_\pm$ and we have
    $T(P_\pm)=P_\pm$.
\end{proof}

For a map $T$ on $\mcM_d(\CC)$ let us define the set of fixed points
\begin{equation}\tag{6.48}
    \FT := \{X\in\mcM_d(\CC)\mid X=T(X)\}.
\end{equation}
Note that if $T$ is linear, then $\FT$ is closed under linear combination, i.e.,
it is a vector space, and if $T$ is Hermiticity preserving (in particular, if it
is positive), then $\FT$ is closed under Hermitian conjugation. The previous
Prop.~6.8 immediately implies that for trace-preserving quantum channels the set
of fixed points is spanned by stationary density operators:

\begin{corollary}[Linearly independent stationary states]
    Let $T:\mcM_d(\CC)\to\mcM_d(\CC)$ be a trace-preserving, positive, linear map. If
    the space $\FT$ has dimension $D$, then there are $D$ linearly independent
    density operators $\{\rho_\alpha\in\mcM_d(\CC)\}_{\alpha=1,\dots,D}$ such that
    $T(\rho_\alpha)=\rho_\alpha$ and
    \begin{equation}\tag{6.49}
        \FT = \spn\{\rho_1,\dots,\rho_D\}.
    \end{equation}
\end{corollary}

\begin{proposition}[Maximal support of fixed points]
    Let $T:\mcM_d(\CC)\to\mcM_d(\CC)$ be a trace-preserving, positive, linear map and
    $T_\infty$ the projection onto its fixed point space as given in Eq.~\eqref{6.14}.
    Then every fixed point $X=T(X)\in\mcM_d(\CC)$ has support and range within the
    support space of the fixed point $T_\infty(\one)$.
\end{proposition}

\begin{proof}
    For positive $X$ the claim follows from positivity of $T_\infty$: since
    $0\le T_\infty(X)=X\le \lVert X\rVert_\infty T_\infty(\one)$ the range (which
    equals the support) of $X$ has to be contained in the one of $T_\infty(\one)$
    (see Douglas' theorem~5.1).

    If $X$ is not positive we can by Prop.~6.8 decompose it into a linear combination
    of four positive fixed points $P_1,\dots,P_4$. Using that $\supp(X)$ is in the
    linear span of $\bigcup_j \supp(P_j)$ together with the result for positive fixed
    points then shows $\supp(X)\subseteq \supp(T_\infty(\one))$. The same argument
    applies for the range of $X$.

    Alternatively, we may use that $X$ is a fixed point iff $X^\dagger$ is, and that
    the range of $X$ equals the support of $X^\dagger$.
\end{proof}

The following shows that the ranges of all stationary states play an exceptional
role: every state which is supported within such a space will remain so after
the action of the channel:

\begin{proposition}[Stationary subspaces]
    Let $T:\mcM_d(\CC)\to\mcM_d(\CC)$ be a trace-preserving, positive, linear map, and
    let $Q\in\mcM_d(\CC)$ be the Hermitian projection onto the support of any
    stationary density operator $\sigma=T(\sigma)\in\mcM_d(\CC)$. Then
    $\tr[(\one-Q)T(Q)]=0$ and for every density operator $\rho\in\mcM_d(\CC)$
    \begin{equation}\tag{6.50}
        \rho\le Q \ \Rightarrow\ T(\rho)\le Q.
    \end{equation}
\end{proposition}

\begin{proof}
    By definition of $Q$ we have that $Q\le c\sigma\le c'Q$ for some strictly positive
    numbers $c$ and $c'$. Therefore $T(Q)\le cT(\sigma)=c\sigma\le c'Q$ which implies
    $\tr[(\one-Q)T(Q)]=0$. Similarly, if $\rho\le Q$, then $T(\rho)\le c'Q$, i.e.,
    $T(\rho)$ has support in the range of $Q$. Since $T(\rho)$ has eigenvalues at most
    one, this means that $T(\rho)\le Q$.
\end{proof}

Projections onto stationary subspaces (in the sense of Eq.~\eqref{6.50}) can be
easily characterized in the Heisenberg picture, as well:

\begin{proposition}[Stationary subspaces II]
    Let $T:\mcM_d(\CC)\to\mcM_d(\CC)$ be a trace-preserving, positive, linear map, and
    let $Q\in\mcM_d(\CC)$ be any Hermitian projection. Then the following statements
    are equivalent:
    \begin{enumerate}
        \item For every density matrix $\rho\in\mcM_d(\CC)$ with $\rho\le Q$, we have that
              $T(\rho)\le Q$ as well,
        \item $T^*(Q)\ge Q$.
    \end{enumerate}
\end{proposition}

\begin{proof}
    $1.\Rightarrow 2$.: define the orthogonal projection $P:=(\one-Q)$. The assumption
    implies that $Q\,T^*(P)\,Q=0$ which in turn means that $P\,T^*(P)\,P=T^*(P)$. Using
    that $T^*(\one)=\one$, the latter can be rewritten as $T^*(Q)-P\,T^*(Q)\,P=Q$ from
    which $T^*(Q)\ge Q$ follows.
    $2.\Rightarrow 1$. follows from the chain of inequalities
    \[
        0 \le \tr[T(\rho)P]
        = 1-\tr[\rho\,T^*(Q)]
        \le 1-\tr[\rho Q]
        = 0.
    \]
\end{proof}

A projection onto a stationary subspace is often called \emph{sub-harmonic} and
item~2.\ of the previous proposition is used as the defining property of this
notion. Note that while the support projection of every stationary state is
sub-harmonic (due to Props.~6.10,6.11), the converse is not necessarily true:
for instance $T^*(\one)\ge \one$ holds trivially for every trace-preserving $T$
even if there is no stationary state of full rank.

If $T$ is completely positive with Kraus decomposition
$T(\cdot)=\sum_i K_i\,\cdot\,K_i^\dagger$, one can relate support spaces of fixed
points to properties of the Kraus operators: suppose there is a stationary state
with support $\mcH\subset \CC^d$. Then all Kraus operators have to preserve this
subspace in the sense that $K_i\mcH\subseteq \mcH$. In particular, if
$\rho=\ket{\psi}\!\bra{\psi}$ is a pure state, then $\psi$ has to be a
simultaneous eigenvector of all Kraus operators.

We will now investigate the structure of the set of fixed points $\FT$ of
quantum channels. As a first step we exploit Props.~6.9,6.10 in order to restrict
the discussion to trace-preserving maps which have a full rank fixed point:

\begin{lemma}[Restriction to full rank fixed points]
    Let $T:\mcM_d(\CC)\to\mcM_d(\CC)$ be a trace-preserving, (completely) positive,
    linear map, and $\mcH\subseteq \CC^d$ the support space of $T_\infty(\one)$. Then
    the linear map $\widetilde{T}:\mcB(\mcH)\to\mcB(\mcH)$ defined via
    \[
        \widetilde{T}(X) := \left.T(X\oplus 0)\right|_{\mcH}
    \]
    is trace-preserving and (completely) positive as well. Moreover, $\widetilde{T}$
    has a fixed point of full rank and\footnote{Here the `0' is supposed to act on
        the orthogonal complement of $\mcH$ in $\CC^d$.}
    \begin{equation}\tag{6.51}
        \FT = \mathcal{F}_{\widetilde{T}} \oplus 0.
    \end{equation}
\end{lemma}

\begin{proof}
    Denote by $V:\mcH\to\CC^d$ the isometry which embeds $\mcH$ into $\CC^d$ so that
    $V^\dagger V=\one$ and $Q:=VV^\dagger$ is the projection onto the range of
    $T_\infty(\one)$. Then by definition $\widetilde{T}(X)=V^\dagger T(VXV^\dagger)V$
    so that $\widetilde{T}$ inherits the property of being (completely) positive from
    $T$. In order to see that $\widetilde{T}$ is trace-preserving note that
    $VXV^\dagger=X\oplus 0$ has support and range within the range of $Q$. Thus we
    can decompose it into four positive operators with the same property and apply
    Prop.~6.10 which leads to
    \begin{equation}\tag{6.52}
        T(X\oplus 0) = \widetilde{T}(X)\oplus 0,
    \end{equation}
    so that the trace-preserving property is inherited from $T$ as well. Finally,
    Eq.~\eqref{6.52} also implies Eq.~\eqref{6.51} since by Prop.~6.9 all fixed points
    of $T$ are of the form $X\oplus 0$. In particular, by definition of $\mcH$, there
    is one fixed point which has full rank on $\mcH$.
\end{proof}

Before returning to trace-preserving maps we will discuss their duals---unital
maps. This will provide an important prerequisite for the discussion of the
trace-preserving case. We will also see that the structure of the fixed point set
of a map simplifies a lot when the dual map has at least one fixed point of full
rank. This is already visible in the dual analog of Lem.~6.4:

\begin{lemma}
    Let $T:\mcM_d(\CC)\to\mcM_d(\CC)$ be a trace-preserving, positive, linear map, and
    $\widetilde{T}$ and $V$ defined as in Lem.~6.4. Let us further denote by $S$ the
    set of linear unital maps which satisfy the Schwarz inequality in Eq.~(5.2). Then
    \begin{equation}\tag{6.53}
        A\in \mathcal{F}_{T^*}\ \Rightarrow\ V^\dagger A V \in \mathcal{F}_{\widetilde{T}^*},
    \end{equation}
    and
    \begin{equation}\tag{6.54}
        X\in \mathcal{F}_{\widetilde{T}^*}
        \ \Leftrightarrow\
        V^\dagger T^*(VXV^\dagger)V = X.
    \end{equation}
    Moreover,
    \begin{equation}\tag{6.55}
        T^*\in S\ \Rightarrow\ \widetilde{T}^*\in S.
    \end{equation}
\end{lemma}

\begin{proof}
    Eq.~\eqref{6.54} follows from inserting the definitions. For Eq.~\eqref{6.53}
    note that
    $\widetilde{T}^*(V^\dagger A V)=V^\dagger T^*(QAQ)V$ (using the notation of
    Lem.~6.4 again). Taking the trace with an arbitrary $C\in\mcB(\mcH)$ we obtain
    \begin{align}
        \tr\![C\,\widetilde{T}^*(V^\dagger A V)]
         &= \tr\![T(VCV^\dagger)\,QAQ], \tag{6.56}\\
         &= \tr\![T(VCV^\dagger)\,A]
        = \tr\![CV^\dagger A V], \tag{6.57}
    \end{align}
    where the second equality follows from Prop.~6.10 and the last step uses the
    assumption that $A$ is a fixed point of $T^*$.

    Finally, the Schwarz inequality for $\widetilde{T}^*$ follows from
    \begin{align}
        V^\dagger T^*(VA^\dagger AV^\dagger)V
         &\ge V^\dagger T^*(VA^\dagger V^\dagger)\,T^*(VAV^\dagger)V, \tag{6.58}\\
         &\ge V^\dagger T^*(VA^\dagger V^\dagger)V
        V^\dagger T^*(VAV^\dagger)V. \tag{6.59}
    \end{align}
    Here the first inequality comes from first inserting $V^\dagger V=\one$ between
    $A^\dagger$ and $A$ and then using the Schwarz inequality for $T^*$ and the
    second inequality is due to $VV^\dagger\le \one$. Written in terms of
    $\widetilde{T}^*$ the above equations lead to the claimed
    $\widetilde{T}^*(A^\dagger A)\ge \widetilde{T}^*(A^\dagger)\widetilde{T}^*(A)$.
\end{proof}

Due to the aforementioned issue we will first characterize the fixed point set
of unital maps which have a dual with fixed point of full rank:

\begin{theorem}[Structure of fixed points for unital Schwarz maps]
    Let $T^*:\mcM_d(\CC)\to\mcM_d(\CC)$ be a unital map which (i) satisfies the
    Schwarz inequality Eq.~(5.2) (e.g., a completely positive map) and (ii) is such
    that $T$ has a full rank fixed point. Then the set of fixed points $\mathcal{F}_{T^*}$ is
    a $*$-algebra and has thus the form specified in Eq.~(1.39) (without the
    zero-block).
\end{theorem}

\begin{proof}
    Let $\rho\in\mcM_d(\CC)$ be a full rank fixed point of $T$. Note that by
    Prop.~6.9 we can w.l.o.g.\ chose a positive definite $\rho>0$. Denote by
    $A=T^*(A)\in\mcM_d(\CC)$ a fixed point of $T^*$. Then, using the fixed point
    properties and the fact that $A^\dagger$ is a fixed point of $T^*$ as well we obtain
    \begin{equation}\tag{6.60}
        \tr\!\left[(T^*(A^\dagger A)-T^*(A^\dagger)T^*(A))\rho\right]=0.
    \end{equation}
    Note that the expression in parentheses is positive semi-definite by the Schwarz
    inequality Eq.~(5.2) and therefore, in fact, identical zero since $\rho$ is assumed
    to be positive definite. Hence, we have equality in the Schwarz inequality so that
    we can apply Thm.~5.4 which implies that if $A,X\in \mathcal{F}_{T^*}$, then $AX\in \mathcal{F}_{T^*}$
    is a fixed point as well. Since $\mathcal{F}_{T^*}$ is also closed under linear combinations
    and Hermitian conjugation, it is a $*$-algebra. Evidently, it is unital since
    $T^*(\one)=\one$ holds by assumption.
\end{proof}

If $T$ does not have a full rank fixed point, then $\mathcal{F}_{T^*}$ need not be an
algebra. Together with Lem.~6.5 we obtain a weaker structure in this case:

\begin{corollary}
    Let $T^*:\mcM_d(\CC)\to\mcM_d(\CC)$ be a unital linear map which satisfies the
    Schwarz inequality Eq.~(5.2). Let $Q=Q^\dagger\in\mcM_d(\CC)$ be the projection
    onto the maximum rank fixed point of $T$, then the following set is a
    $*$-algebra:
    \begin{equation}\tag{6.61}
        \bigl\{Y\in Q\mcM_d(\CC)Q \big| QT^*(Y)Q=Y\bigr\}.
    \end{equation}
\end{corollary}

\begin{proof}
    The stated set is exactly the one appearing in Eq.~\eqref{6.54}, i.e., the fixed
    point set of $\widetilde{T}^*$. Since $\widetilde{T}$ has, however, a fixed point
    of full rank by construction (see Lem.~6.4) we can apply Thm.~6.12 which proves
    the claim.
\end{proof}

The relation between fixed point sets of unital maps and $*$-algebras becomes
more transparent for completely positive maps:

\begin{theorem}[Fixed points and commutants]
    Let $T^*:\mcM_d(\CC)\to\mcM_d(\CC)$ be a completely positive unital map with Kraus
    decomposition $T^*(\cdot)=\sum_i K_i^\dagger\,\cdot\,K_i$. If both $X,X^\dagger X\in
        \mathcal{F}_{T^*}$, then $[X,K_i]=0$ for all $i$. Consequently, the largest $*$-subalgebra
    contained within $\mathcal{F}_{T^*}$ is given by the set
    \begin{equation}\tag{6.62}
        \bigl\{X\in\mcM_d(\CC) \big| \forall i:\ [X,K_i]=[X,K_i^\dagger]=0\bigr\}.
    \end{equation}
\end{theorem}

\begin{proof}
    The assertion follows from the equation
    \[
        \sum_i [X,K_i]^\dagger [X,K_i]
        = T^*(X^\dagger X) - X^\dagger T^*(X) - T^*(X)^\dagger X + X^\dagger X
        = 0
    \]
    by exploiting the fixed point properties of $X$ and $X^\dagger X$ together with
    the fact that the l.h.s.\ is a sum over positive terms which is zero iff all
    terms vanish individually. Hence every $*$-subalgebra of $\mathcal{F}_{T^*}$ is a subset
    of Eq.~\eqref{6.62}. Conversely, the set in Eq.~\eqref{6.62} is a $*$-subalgebra
    within $\mathcal{F}_{T^*}$ on its own and thus the largest one.
\end{proof}

Finally, we discuss the structure of the fixed point set of trace-preserving maps:

\begin{theorem}[Fixed points of trace-preserving maps]
    Let $T:\mcM_d(\CC)\to\mcM_d(\CC)$ be a trace-preserving, positive, linear map for
    which $T^*$ satisfies the Schwarz inequality in Eq.~(5.2). Then there is a
    unitary $U\in\mcM_d(\CC)$ and a set of positive definite density matrices
    $\rho_k\in \mcM_{m_k}(\CC)$ such that the fixed point set of $T$ is given by
    \begin{equation}\tag{6.63}
        \FT
        = U\left(0 \oplus \bigoplus_{k=1}^K \mcM_{d_k}\otimes \rho_k\right)U^\dagger,
    \end{equation}
    for an appropriate decomposition of the Hilbert space
    $\CC^d = \CC^{d_0}\oplus \bigoplus_k \CC^{d_k}\otimes \CC^{m_k}$. $\mcM_{d_k}$
    stands for the full algebra of complex matrices on $\CC^{d_k}$.
\end{theorem}

\begin{proof}
    We exploit that by Lem.~6.4 $\FT = 0\oplus \mathcal{F}_{\widetilde{T}}$ after an
    appropriate basis transformation which we can assign to a unitary $U$. Now
    $\widetilde{T}$ has, by construction, a full rank fixed point and following
    Lem.~6.4 it is trace-preserving and its dual satisfies the Schwarz inequality
    due to Lem.~6.5. This implies by Thm.~6.12 that $\widetilde{T}_\infty^*$ projects
    onto a $*$-algebra. The general form of a positive map projecting onto a
    $*$-algebra (i.e., a conditional expectation) is given in Eq.~(1.40). Taking the
    dual of this map we arrive at $\widetilde{T}_\infty$ projecting onto a set of the
    form in Eq.~\eqref{6.63} (without the zero-block). Putting things together and
    using that we can always move the kernel of each $\rho_k$ into the zero-block
    completes the proof.
\end{proof}

Note that Thm.~6.14 implies that the fixed point $T_\infty(\one)$ is (up to
normalization) given by
$U(0 \oplus \bigoplus_k \one_{m_k}\otimes \rho_k)U^\dagger$. Moreover,
it shows that again the set of fixed points is closely related to a $*$-algebra
(actually is a $*$-algebra w.r.t.\ a modified product):

\begin{corollary}
    Let $T:\mcM_d(\CC)\to\mcM_d(\CC)$ be a trace-preserving, positive, linear map for
    which $T^*$ satisfies the Schwarz inequality in Eq.~(5.2). Then for any density
    matrix $\rho\in\mcM_d(\CC)$ which is a maximum rank fixed point of $T$ we have
    that the set $\rho^{-1/2}\FT\rho^{-1/2}$ is a $*$-algebra (with the inverse taken
    on the support of $\rho$).
\end{corollary}

Another interesting consequence of Thm.~6.14 is its implication according to
Prop.~6.10: the decomposition of the fixed point space into a direct sum in
Eq.~\eqref{6.63} implies the same decomposition into stationary subspaces. What
is more, the $\mcM_{d_k}$'s give rise to subsystem which are invariant, i.e., on
which $T$ acts as the identity (up to the reversible rotation given by $U$).

\begin{theorem}[Unique fixed points of full rank]
    Let $T:\mcM_d(\CC)\to\mcM_d(\CC)$ be a trace-preserving quantum channel with Kraus
    decomposition $T(\cdot)=\sum_i K_i\,\cdot\,K_i^\dagger$. If for some $n\in\NN$ we
    have that
    \begin{equation}\tag{6.64}
        \spn\left\{\prod_{k=1}^n K_{i_k}\right\}=\mcM_d(\CC),
    \end{equation}
    i.e., homogeneous polynomials of the Kraus operators span the entire matrix
    algebra, then there exist a unique positive definite density matrix
    $\rho\in\mcM_d(\CC)$ such that $\FT\propto \rho$.
\end{theorem}

\begin{proof}
    Consider the map $T^n$, i.e., the $n$-fold concatenation of $T$. This has Kraus
    operators $R_I:=\prod_{k=1}^n K_{i_k}$ with multi-index $I=(i_1,\dots,i_n)$ and it
    satisfies $\FT\subseteq \mathcal{F}_{T^n}$. Now for every $\psi\in\CC^d$ we have that
    $T^n(\ket{\psi}\!\bra{\psi})$ has full rank since by assumption
    $\spn\{R_I\ket{\psi}\}=\CC^d$. This implies that there cannot be any stationary
    density operator which has a kernel. Since we can always construct a rank deficient
    positive fixed point if the dimension of $\FT$ is larger than one (see Cor.~6.5),
    the assertion follows.
\end{proof}

\begin{problem}[Fixed points of general positive maps]
Characterize the set of fixed points of unital, positive, linear maps which do
not satisfy the Schwarz inequality.
\end{problem}

% ------------------------------------------------------------
\subsection{Cycles and recurrences}
% ------------------------------------------------------------

We will now have a closer look at the space corresponding to the peripheral
spectrum (i.e., eigenvalues which are phases) of a positive linear map
$T:\mcM_d(\CC)\to\mcM_d(\CC)$ which is assumed to be either trace-preserving or
unital. Denote the complex linear span of all the respective eigenspaces by
\begin{equation}\tag{6.65}
    \XT := \spn\left\{X\in\mcM_d(\CC) \Big| \exists \varphi\in\RR:\ T(X)=e^{i\varphi}X\right\},
\end{equation}
for which we will simply write $\mcX$ if the dependence on $T$ is clear from the
context. The questions which we address first are: what's the structure of
$\mcX$? what's the structure of $T$ restricted to $\mcX$? and what does this
imply on the peripheral spectrum?

Note that evidently $\XT=\mcX_{T_\varphi}$ if $T_\varphi$ is constructed from $T$
as defined in Eq.~\eqref{6.13} by discarding all but those terms in the spectral
decomposition which correspond to the peripheral spectrum. Recall that by
Prop.~6.3 if $T$ is trace-preserving or unital, then $T_\varphi$ will be
(completely) positive if $T$ is and that by Prop.~6.2 there cannot any non-trivial
Jordan-block associated to an eigenvalue of modulus one.

\begin{proposition}[Asymptotic image]
    Let $T:\mcM_d(\CC)\to\mcM_d(\CC)$ be a positive and trace-preserving linear map.
    Then $\XT$ is the image of the projection $T_\phi$, it is spanned by positive
    operators, and it is closed under the action of $T$. That is,
    \begin{enumerate}
        \item $T_\phi(\mcM_d(\CC))=\XT$,
        \item $\exists\{\rho_i\ge 0\}$ such that $\XT=\spn\{\rho_i\}$,
        \item $T(\XT)=\XT$.
    \end{enumerate}
\end{proposition}

\begin{proof}
    Note that by the proof of Prop.~6.3 every $X\in \XT$ is a fixed point of $T_\phi$
    since there is an ascending sequence $n_i\in\NN$ such that
    $\lim_{i\to\infty} T^{n_i}=T_\phi$. Conversely, if $X=T_\phi(X)$ is such a fixed
    point, then it has to be a linear combination of eigenvectors of $T$ which
    correspond to eigenvalues of modulus one. $\XT$ thus coincides with the space of
    fixed points of $T_\phi$ which is in turn given by $T_\phi(\mcM_d(\CC))$, proving 1.

    2.\ follows from 1.\ together with Cor.~6.5 and in order to arrive at 3.\ we just
    have to use the definition of $\XT$ and express an element $X\in\XT$ as a linear
    combination of peripheral eigenvectors of $T$.
\end{proof}

The following refines this Proposition by adding the assumption that the map
under consideration satisfies the Schwarz inequality. Exploiting this fact one
observes that $\mcX$ is a $*$-algebra (w.r.t.\ a modified product) and that $T$
acts on it either by permuting blocks or by unitary conjugation on subsystems:

\begin{theorem}[Structure of cycles]
    Let $T:\mcM_d(\CC)\to\mcM_d(\CC)$ be a trace-preserving and positive linear map
    which fulfills the Schwarz inequality in Eq.~(5.2).
    \begin{enumerate}
        \item There exists a decomposition of the Hilbert space
              $\CC^d=\mcH_0\oplus \bigoplus_{k=1}^K \mcH_k$ into a direct sum of tensor products
              $\mcH_k=\mcH_{k,1}\otimes \mcH_{k,2}$ and positive definite density matrices
              $\rho_k$ acting on $\mcH_{k,2}$ such that
              \begin{equation}\tag{6.66}
                  \XT = 0 \oplus \bigoplus_{k=1}^K \mcM_{d_k}\otimes \rho_k,
              \end{equation}
              where $\mcM_{d_k}$ is a full complex matrix algebra on $\mcH_{k,1}$ of dimension
              $d_k:=\dim(\mcH_{k,1})$. That is, for every $X\in\XT$ there are
              $x_k\in\mcB(\mcH_{k,1})$ such that
              \begin{equation}\tag{6.67}
                  X = 0 \oplus \bigoplus_k x_k\otimes \rho_k.
              \end{equation}
        \item There exist unitaries $U_k\in\mcB(\mcH_{k,1})$ and a permutation $\pi$,
              which permutes within subsets of $\{1,\dots,K\}$ for which the corresponding
              $\mcH_k$'s have equal dimension, so that for every $X\in\XT$ represented as in
              Eq.~\eqref{6.67}
              \begin{equation}\tag{6.68}
                  T(X)
                  = 0 \oplus \bigoplus_{k=1}^K
                  U_k\,x_{\pi(k)}\,U_k^\dagger \otimes \rho_k.
              \end{equation}
    \end{enumerate}
\end{theorem}

\begin{proof}
    The basic ingredients of the proof are those which already appeared in the proof
    of Thm.~6.1 where we characterized maps with unit determinant. Exploiting
    Dirichlet's lemma on Diophantine approximations (Lem.~6.1) we can find an
    ascending subsequence $n_i\in\NN$ so that $\lim_{i\to\infty} T^{n_i}$ converges
    to a map $I$ with eigenvalues which are either one or zero. Clearly, $I$ is again
    a trace-preserving, positive linear map which satisfies the Schwarz inequality
    since all these properties are preserved under concatenation. Moreover,
    $\XT=\mcX_I$ is the fixed-point space of $I$. So, statement~1.\ of the theorem
    follows from the structure of the space of fixed points provided in Thm.~6.14.

    2.\ The map $T^{-1}:=\lim_{i\to\infty} T^{n_i-1}$ is the inverse of $T$ on $\mcX$
    since by construction $T^{-1}T=I$. Hence, both $T$ and $T^{-1}$ map $\mcX\to\mcX$
    in a bijective way. The crucial point here is that $T^{-1}$ is again a positive,
    trace-preserving map since it is constructed as a limit of such maps. To
    understand the consequences, consider any pure state in $\mcX$, i.e., a density
    operator $\sigma\in\mcX$ which has no non-trivial convex decomposition within
    $\mcX$. Then the image of $\sigma$ under $T$ has to be a pure state as well:
    assume this is not the case, i.e., $T(\sigma)=\sum_i \lambda_i \sigma_i$ is a
    non-trivial convex decomposition into states $\sigma_i\in\mcX$. Then applying
    $T^{-1}$ to this equation leads to a contradiction since
    $\sigma=\sum_i \lambda_i T^{-1}(\sigma_i)$ is not pure. Consequently, both $T$
    and $T^{-1}$ map pure states in $\mcX$ onto pure states. Note that a pure state
    can only have support in one of the $K$ blocks, for instance $\sigma=x\otimes\rho_k$
    where $x\in\mcB(\mcH_{k,1})$ is a rank one projection. Now we know that
    $T(\sigma)=x'\otimes \rho_{k'}$ for some $k'$ and some rank one projection
    $x'\in\mcB(\mcH_{k',1})$. By continuity and the fact that $T$ is a bijective linear
    map on $\mcX$, we have that within $\mcX$: (i) every element of block $k$ is mapped
    to an element of the block $k'$ (i.e., $k'$ does not depend on $x$), and (ii) the
    blocks $k$ and $k'$ must have equal dimension $d_k=d_{k'}$. Therefore, there is a
    permutation $\pi$ which permutes blocks of equal size so that $X\in\mcX$ is mapped to
    \[
        T(X) = 0 \oplus \bigoplus_k T_k(x_{\pi(k)})\otimes \rho_k,
    \]
    with some linear maps $T_k:\mcM_{d_k}\to\mcM_{d_k}$. Since the latter have, together
    with their inverses, to be positive and trace-preserving they must be either matrix
    transpositions or unitary conjugations by Cor.~6.2. Matrix transpositions are,
    however, ruled out by the requirement that $T$ and thus each $T_k$ is a Schwarz map.
\end{proof}

% ------------------------------------------------------------
\subsection{Inverse eigenvalue problems}
% ------------------------------------------------------------

% (No additional content in the provided pages.)

% ------------------------------------------------------------
\subsection{Literature}
% ------------------------------------------------------------

There are various extensions of Brouwer's fixed point theorem. Schauder's theorem,
for instance, extends it to topological vector spaces, which may be of infinite
dimension. That is, $S$ becomes a convex, compact subset of a topological vector
space. Kakutani's theorem is a generalization to upper semi-continuous set-valued
functions (correspondences) $T$ from a non-empty, compact, convex set $S\subset\RR^n$
into itself such that for all $x\in S$ the set $T(x)$ is non-empty and convex, too.
Dealing with these and more refined fixed point theorems is part of algebraic topology.

\end{document}
